\documentclass[a4paper]{article} \pagestyle{plain}
\setlength\parindent{0ex}
\usepackage{url}
\usepackage{graphicx}
\usepackage{amssymb,amsmath}






\renewcommand{\maketitle}{

\begin{titlepage}

\begin{center}
\includegraphics[width=14cm]{../../_static/logos}
\end{center}

\sf{

\vspace*{2.0cm}
\noindent{
\huge{
\textbf{
The `Case-By-Case' Schema for Molecular States in XSAMS - v0.2
}}}
\vspace*{3.0cm}

}

\noindent \textbf{Document Information}

\vspace*{0.5cm}

\setlength{\parindent}{-1cm}

\noindent \begin{tabular}{p{1.7in}p{4.3in}} %\hline 
\textbf{Editors: } & C. Hill \\
\textbf{Authors:} & C. Hill \\ %\hline 
\textbf{Contributors:} & M.-L. Dubernet, J. Tennyson, E. Roueff \\ %\hline 
\textbf{Type of document:} & standards documentation \\ %\hline 
\textbf{Status:} & draft \\ %\hline 
\textbf{Distribution:} & public \\ %\hline 
\textbf{Work package:} & WP6 \\ %\hline 
\textbf{Version:} & 0.2 \\ %\hline 
\textbf{Document code:} &  \\ %\hline 
\textbf{Directory and file name:} & \url{http://www.vamdc.org/documents/cbc_v0.2.pdf}  \\ %\hline 

\end{tabular}

\vspace*{2.0cm}

\noindent \begin{tabular}{p{1.7in}p{4.3in}}

\textbf{Abstract:} & 
This document describes the `case-by-case' XML Schema for molecular states within VAMDC-XSAMS.
\end{tabular}



\end{titlepage}

\noindent \textbf{Version History}

\textbf{}

\noindent \begin{tabular}{|l|l|l|l|} 
\hline 
\textbf{Version} & \textbf{Date} & \textbf{Modified By} & \textbf{Description of Change} \\ \hline 
0.1 & 16/12/2010 & C. Hill & first draft \\ \hline 
0.2 & 03/05/2011 & C. Hill & second draft \\ \hline
\end{tabular}

\textbf{}

\noindent \textbf{Disclaimer}

\noindent The information in this document is subject to change without notice. Company or product names mentioned in this document may be trademarks or registered trademarks of their respective companies.


\textbf{}

\noindent \textbf{All rights reserved}

\noindent The document is proprietary of the VAMDC consortium members. No copying or distributing, in any form or by any means, is allowed without the prior written agreement of the owner of the property rights.

\noindent 

\noindent This document reflects only the authors' view. The European Community is not liable for any use that may be made of the information contained herein.

\textbf{}

\noindent \textbf{Acknowledgements}

\noindent VAMDC is funded under the ``Combination of Collaborative Projects and Coordination and  Support Actions'' Funding Scheme of The Seventh Framework Program. Call topic: INFRA-2008-1.2.2 Scientific Data Infrastructure. Grant Agreement number: 239108.

\textbf{}

}









\begin{document}

\maketitle
\pagebreak

\section*{Introduction}

The `case-by-case' XML description of molecular states within XSAMS is designed to provide a straightforward and `flat' data structure for representing the quantum numbers and symmetries that define a molecular state. As of version 0.2.1, 12 cases have been implemented, covering the needs of the HITRAN, CDMS, and BASECOL databases.

Each case is identified by a \texttt{prefix} and a \texttt{version}, and belongs to the namespace given (currently) by the URI \url{http://www.ucl.ac.uk/~ucapch0/XSAMS/cases/<prefix>/<version>}. At least for the time being, it is suggeested that validation is through Namespace Validation Dispatch Language (NVDL). This may be implemented by including the relevant processing instruction in the XSAMS XML instance document, after the XML declaration. For example, using the oxygen editor
\begin{verbatim}
<?xml version="1.0" encoding="UTF-8"?>
<?oxygen NVDLSchema="cbc.nvdl"?>
...
\end{verbatim}
The NVDL Schema document, \texttt{cbc.nvdl} contains a set of rules which link the namespaces encountered in the XML document with the Schemata required to validate them. An sample NVDL document may be downloaded from \url{http://www.ucl.ac.uk/~ucapch0/XSAMS/}.

\section*{The Cases}

The identified cases are described in the following section; some examples are given below. It should be noted that the rovibronic states of different electronic states of a molecule may be described using different cases. For example, the ground, $X^2\Pi$ electronic state of $\mathrm{NO}$ may be described within the \texttt{hunda} case whereas the excited, $A^2\Sigma^+$ electronic state would be better described using the \texttt{hundb} case. The electronic state is identified by its single-letter spectroscopic symbol ($X$, $A$, $a$, $B$, etc.)

There follows a list of the cases identified in version 0.2.1 of XSAMS with some examples:

\begin{enumerate}
\item[0.] General case for arbitrary quantum numbers and symmetry (\texttt{gen})
\item Diatomic closed shell (\texttt{dcs}): CO, $\mathrm{N_2}$, $\mathrm{NO^+}$
\item Hund's case (a) diatomics (\texttt{hunda}): NO, OH [for low $J$]
\item Hund's case (b) diatomics (\texttt{hundb}): $\mathrm{O_2}$, OH [for high $J$]
\item Closed-shell, linear triatomic molecules (\texttt{ltcs}): $\mathrm{CO_2}$, HCN
\item Closed-shell, non-linear triatomic molecules (\texttt{nltcs}): $\mathrm{H_2O}$
\item Closed-shell, symmetric top molecules (\texttt{stcs}): $\mathrm{NH_3}$, $\mathrm{CH_3Cl}$
\item Closed-shell, linear, polyatomic molecules (\texttt{lpcs}): $\mathrm{C_2H_2}$
\item Closed-shell, asymmetric top molecules (\texttt{asymcs}): $\mathrm{C_2H_4}$
\item Open-shell, asymmetric top molecules (\texttt{asymos}): $\mathrm{CH_3O}$
\item Closed-shell, spherical top molecules (\texttt{sphcs}): $\mathrm{CH_4}$, $\mathrm{SF_6}$
\item Open-shell, spherical top molecules (\texttt{sphos})
\item Open-shell, linear triatomic molecules (\texttt{ltos}): $\mathrm{CCH}$
\item Open-shell, linear, polyatomic molecules (\texttt{lpos}): $\mathrm{C_3H}$, $\mathrm{C_{10}H}$
\item Open-shell, non-linear triatomic molecules (\texttt{nltos}): $\mathrm{HO_2}$, $\mathrm{CH_2}$
\end{enumerate}

\section*{Examples}
The element \texttt{<case:QNs>} should be placed within the XSAMS element \texttt{MolecularState} (although this cannot be enforced using NVDL). Some examples follow.

\subsection*{Example 1: a rovibrational state of CO}

The $v=0, J=1$ state of the ground electronic state of CO has the following representation in the `case-by-case' formulism:

\begin{verbatim}
<dcs:QNs>
   <dcs:ElecStateLabel>X</dcs:ElecStateLabel>
   <dcs:v>1</dcs:v>
   <dcs:J>0</dcs:J>
</dcs:QNs>
\end{verbatim}

\subsection*{Example 2: a rovibrational state of $\mathrm{NH_3}$}

The $J=22, K=10$ rotational state of the $(1,0^+,0^0,2^2)$ vibrational level of $\mathrm{NH_3}$ could be represented by the following XML:

\begin{verbatim}
<stcs:QNs>
   <stcs:ElecStateLabel>X</stcs:ElecStateLabel>
   <stcs:vi mode="1">1</stcs:vi>
   <stcs:vi mode="2">0</stcs:vi>
   <stcs:vi mode="3">0</stcs:vi>
   <stcs:vi mode="4">2</stcs:vi>
   <stcs:li mode="3">0</stcs:li>
   <stcs:li mode="4">2</stcs:li>
   <stcs:vibInv>s</stcs:vibInv>
   <stcs:vibSym>E</stcs:vibSym>
   <stcs:J>20</stcs:J>
   <stcs:K>10</stcs:K>
</stcs:QNs>
\end{verbatim}

\subsection*{Example 3: two states of different electronic states of MgH}

This example shows two states of the MgH radical, using different cases for the $X^2\Sigma^+$ and $A^2\Pi$ electronic states:

\begin{verbatim}
<MolecularState stateID="S1-MgH-1">
  <Description>
    A state in the ground electronic state, X(2Sigma+), of MgH
  </Description>
  <MolecularStateCharacterisation>
    <StateEnergy energyOrigin="Zero-point from calculation based on N^2 Hamiltonian">
      <Value units="1/cm">0.</Value>
    </StateEnergy>
    <TotalStatisticalWeight>4</TotalStatisticalWeight>
  </MolecularStateCharacterisation>
  <hundb:QNs>
    <hundb:ElecStateLabel>X</hundb:ElecStateLabel>
    <hundb:Lambda>0</hundb:Lambda>
    <hundb:S>0.5</hundb:S>
    <hundb:v>0</hundb:v>
    <hundb:J>0.5</hundb:J>
    <hundb:N>0</hundb:N>
    <hundb:SpinComponentLabel>1</hundb:SpinComponentLabel>
    <hundb:parity>+</hundb:parity>
    <hundb:kronigParity>e</hundb:kronigParity>
  </hundb:QNs>
</MolecularState>
  ...
<MolecularState stateID="S1001-MgH-1">
  <Description>
    A state in the first excited electronic state, A(2Pi) of MgH
  </Description>
  <MolecularStateCharacterisation>
    <StateEnergy energyOrigin="Zero-point of electronic ground state">
      <Value units="1/cm">19273.2694</Value>
    </StateEnergy>
    <TotalStatisticalWeight>4</TotalStatisticalWeight>
  </MolecularStateCharacterisation>
  <hunda:QNs>
    <hunda:ElecStateLabel>A</hunda:ElecStateLabel>
    <hunda:Lambda>1</hunda:Lambda>
    <hunda:Omega>0.5</hunda:Omega>
    <hunda:S>0.5</hunda:S>
    <hunda:v>0</hunda:v>
    <hunda:J>0.5</hunda:J>
    <hunda:parity>-</hunda:parity>
    <hunda:kronigParity>f</hunda:kronigParity>
  </hunda:QNs>
</MolecularState>
\end{verbatim}

% case-by-case descriptions
%%%

\section*{Diatomic closed-shell molecules: dcs}
\fbox{
\parbox{12cm}{
\subsection*{\texttt{ElecStateLabel}}
\subsubsection*{\underline{Element}}
\texttt{<dcs:ElecStateLabel>}
\subsubsection*{\underline{Attributes}}
None.
\subsubsection*{\underline{Description}}
\texttt{ElecStateLabel} is a label identifying the electronic state: $X$, $A$, $a$, $B$, etc..
\subsubsection*{\underline{Restrictions}}
\begin{itemize}
\item string.
\end{itemize}
}
}
\vspace{0.5cm}
\fbox{
\parbox{12cm}{
\subsection*{$v$}
\subsubsection*{\underline{Element}}
\texttt{<dcs:v>}
\subsubsection*{\underline{Attributes}}
None.
\subsubsection*{\underline{Description}}
$v$ is the vibrational quantum number.
\subsubsection*{\underline{Restrictions}}
\begin{itemize}
\item non-negative integer.
\end{itemize}
}
}
\vspace{0.5cm}
\fbox{
\parbox{12cm}{
\subsection*{$J$}
\subsubsection*{\underline{Element}}
\texttt{<dcs:J>}
\subsubsection*{\underline{Attributes}}
None.
\subsubsection*{\underline{Description}}
$J$ is the quantum number associated with the total angular momentum excluding nuclear spin, $\boldsymbol{J}$.
\subsubsection*{\underline{Restrictions}}
\begin{itemize}
\item non-negative integer.
\end{itemize}
}
}
\vspace{0.5cm}
\fbox{
\parbox{12cm}{
\subsection*{$F_1$}
\subsubsection*{\underline{Element}}
\texttt{<dcs:F1>}
\subsubsection*{\underline{Attributes}}
\texttt{nuclearSpinRef}: a label identifying the nuclear spin coupled to $\boldsymbol{J}$ to form the intermediate angular momentum.
\subsubsection*{\underline{Description}}
$F_1$ is the intermediate angular momentum quantum number associated with the coupling of the rotational angular momentum and nuclear spin of nucleus 1 where two such couplings are resolved: $\boldsymbol{F_1} = \boldsymbol{J} + \boldsymbol{I_1}$; $F_1$ is often not a good quantum number.
\subsubsection*{\underline{Restrictions}}
\begin{itemize}
\item non-negative integer or half-integer.
\item $|J - I_1| \le F_1 \le J + I_1$.
\end{itemize}
}
}
\vspace{0.5cm}
\fbox{
\parbox{12cm}{
\subsection*{$F$}
\subsubsection*{\underline{Element}}
\texttt{<dcs:F>}
\subsubsection*{\underline{Attributes}}
\texttt{nuclearSpinRef}: a label identifying the nuclear spin coupled to $\boldsymbol{J}$ (or $\boldsymbol{F_1}$) to form the total angular momentum.
\subsubsection*{\underline{Description}}
$F$ is the quantum number associated with the total angular momentum including nuclear spin: $\boldsymbol{F} = \boldsymbol{J} + \boldsymbol{I_1}$ if only one such coupling is resolved, $\boldsymbol{F} = \boldsymbol{F_1} + \boldsymbol{I_2}$ if both couplings are resolved.
\subsubsection*{\underline{Restrictions}}
\begin{itemize}
\item non-negative integer or half-integer.
\end{itemize}
}
}
\vspace{0.5cm}
\fbox{
\parbox{12cm}{
\subsection*{$r$}
\subsubsection*{\underline{Element}}
\texttt{<dcs:r>}
\subsubsection*{\underline{Attributes}}
\texttt{name}: a string identifying this ranking index.
\subsubsection*{\underline{Description}}
$r$ is a named, positive integer label identifying the state if no other good quantum numbers or symmetries are known.
\subsubsection*{\underline{Restrictions}}
\begin{itemize}
\item positive integer.
\end{itemize}
}
}
\vspace{0.5cm}
\fbox{
\parbox{12cm}{
\subsection*{\texttt{parity}}
\subsubsection*{\underline{Element}}
\texttt{<dcs:parity>}
\subsubsection*{\underline{Attributes}}
None.
\subsubsection*{\underline{Description}}
\texttt{parity} is the total parity: the parity of the total molecular wavefunction (excluding nuclear spin) with respect to inversion through the molecular centre of mass of all particles' coordinates in the laboratory coordinate system, the $\hat{E}^*$ operation.
\subsubsection*{\underline{Restrictions}}
\begin{itemize}
\item `$+$' or `$-$'.
\end{itemize}
}
}
\vspace{0.5cm}
\fbox{
\parbox{12cm}{
\subsection*{\texttt{asSym}}
\subsubsection*{\underline{Element}}
\texttt{<dcs:asSym>}
\subsubsection*{\underline{Attributes}}
None.
\subsubsection*{\underline{Description}}
\texttt{asSym} is (for diatomic molecules with a centre of inversion) the symmetry of the rovibronic wavefunction: `a' or `s' such that the total wavefunction including nuclear spin is symmetric or antisymmetric with respect to permutation of the identical nuclei ($\hat{P}_{12}$), according to whether they are bosons or fermions respectively.
\subsubsection*{\underline{Restrictions}}
\begin{itemize}
\item `s' or `a'.
\end{itemize}
}
}
\vspace{0.5cm}




\section*{Hund's case (a) diatomics: hunda}
\fbox{
\parbox{12cm}{
\subsection*{\texttt{ElecStateLabel}}
\subsubsection*{\underline{Element}}
\texttt{<hunda:ElecStateLabel>}
\subsubsection*{\underline{Attributes}}
None.
\subsubsection*{\underline{Description}}
\texttt{ElecStateLabel} is a label identifying the electronic state: $X$, $A$, $a$, $B$, etc..
\subsubsection*{\underline{Restrictions}}
\begin{itemize}
\item string.
\end{itemize}
}
}
\vspace{0.5cm}
\fbox{
\parbox{12cm}{
\subsection*{\texttt{elecInv}}
\subsubsection*{\underline{Element}}
\texttt{<hunda:elecInv>}
\subsubsection*{\underline{Attributes}}
None.
\subsubsection*{\underline{Description}}
\texttt{elecInv} is the parity of the electronic wavefunction with respect to inversion through the molecular centre of mass in the molecular coordinate system.
\subsubsection*{\underline{Restrictions}}
\begin{itemize}
\item `g' or `u'.
\end{itemize}
}
}
\vspace{0.5cm}
\fbox{
\parbox{12cm}{
\subsection*{\texttt{elecRefl}}
\subsubsection*{\underline{Element}}
\texttt{<hunda:elecRefl>}
\subsubsection*{\underline{Attributes}}
None.
\subsubsection*{\underline{Description}}
\texttt{elecRefl} is the parity of the electronic wavefunction with respect to reflection in a plane containing the molecular symmetry axis in the molecular coordinate system (equivalent to inversion through the molecular centre of mass in the laboratory coordinate system).
\subsubsection*{\underline{Restrictions}}
\begin{itemize}
\item `$+$' or `$-$'.
\end{itemize}
}
}
\vspace{0.5cm}
\fbox{
\parbox{12cm}{
\subsection*{$|\mathit{\Lambda}|$}
\subsubsection*{\underline{Element}}
\texttt{<hunda:Lambda>}
\subsubsection*{\underline{Attributes}}
None.
\subsubsection*{\underline{Description}}
$|\mathit{\Lambda}|$ is the quantum number associated with the magnitude of the projection of the total electronic orbital angular momentum, $\boldsymbol{L}$, onto the molecular axis.
\subsubsection*{\underline{Restrictions}}
\begin{itemize}
\item non-negative integer.
\end{itemize}
}
}
\vspace{0.5cm}
\fbox{
\parbox{12cm}{
\subsection*{$|\mathit{\Sigma}|$}
\subsubsection*{\underline{Element}}
\texttt{<hunda:Sigma>}
\subsubsection*{\underline{Attributes}}
None.
\subsubsection*{\underline{Description}}
$|\mathit{\Sigma}|$ is the quantum number associated with the magnitude of the projection of $\boldsymbol{S}$ onto the molecular axis.
\subsubsection*{\underline{Restrictions}}
\begin{itemize}
\item non-negative integer or half-odd integer.
\item $|\mathit{\Sigma}| = S, S-1, \cdots, \frac{1}{2}\;\mathrm{or}\;0$.
\end{itemize}
}
}
\vspace{0.5cm}
\fbox{
\parbox{12cm}{
\subsection*{$\mathit{\Omega}$}
\subsubsection*{\underline{Element}}
\texttt{<hunda:Omega>}
\subsubsection*{\underline{Attributes}}
None.
\subsubsection*{\underline{Description}}
$\mathit{\Omega}$ is the quantum number associated with the projection of the total angular momentum (excluding nuclear spin), $\boldsymbol{J}$, onto the molecular axis: $\mathit{\Omega} = |\mathit{\Lambda} + \mathit{\Sigma}|$ (or $\mathit{\Omega} = |\mathit{\Lambda}| + \mathit{\Sigma}$ if $S > |\mathit{\Lambda}| > 0$).
\subsubsection*{\underline{Restrictions}}
\begin{itemize}
\item non-negative integer or half-integer.
\item $|\mathit{\Omega}| \le J$.
\end{itemize}
}
}
\vspace{0.5cm}
\fbox{
\parbox{12cm}{
\subsection*{$S$}
\subsubsection*{\underline{Element}}
\texttt{<hunda:S>}
\subsubsection*{\underline{Attributes}}
None.
\subsubsection*{\underline{Description}}
$S$ is the quantum number associated with the total electronic spin angular momentum, $\boldsymbol{S}$.
\subsubsection*{\underline{Restrictions}}
\begin{itemize}
\item non-negative integer or half-odd integer.
\end{itemize}
}
}
\vspace{0.5cm}
\fbox{
\parbox{12cm}{
\subsection*{$v$}
\subsubsection*{\underline{Element}}
\texttt{<hunda:v>}
\subsubsection*{\underline{Attributes}}
None.
\subsubsection*{\underline{Description}}
$v$ is the vibrational quantum number.
\subsubsection*{\underline{Restrictions}}
\begin{itemize}
\item non-negative integer.
\end{itemize}
}
}
\vspace{0.5cm}
\fbox{
\parbox{12cm}{
\subsection*{$J$}
\subsubsection*{\underline{Element}}
\texttt{<hunda:J>}
\subsubsection*{\underline{Attributes}}
None.
\subsubsection*{\underline{Description}}
$J$ is the quantum number associated with the total angular momentum excluding nuclear spin: $\boldsymbol{J}=\boldsymbol{L}+\boldsymbol{S}+\boldsymbol{R}$.
\subsubsection*{\underline{Restrictions}}
\begin{itemize}
\item non-negative integer or half-odd integer.
\end{itemize}
}
}
\vspace{0.5cm}
\fbox{
\parbox{12cm}{
\subsection*{$F_1$}
\subsubsection*{\underline{Element}}
\texttt{<hunda:F1>}
\subsubsection*{\underline{Attributes}}
\texttt{nuclearSpinRef}: a label identifying the nuclear spin coupled to $\boldsymbol{J}$ to form the intermediate angular momentum.
\subsubsection*{\underline{Description}}
$F_1$ is the intermediate angular momentum quantum number associated with the coupling of the rotational angular momentum and nuclear spin of nucleus 1 where two such couplings are resolved: $\boldsymbol{F_1} = \boldsymbol{J} + \boldsymbol{I_1}$; $F_1$ is often not a good quantum number.
\subsubsection*{\underline{Restrictions}}
\begin{itemize}
\item non-negative integer or half-integer.
\item $|J - I_1| \le F_1 \le J + I_1$.
\end{itemize}
}
}
\vspace{0.5cm}
\fbox{
\parbox{12cm}{
\subsection*{$F$}
\subsubsection*{\underline{Element}}
\texttt{<hunda:F>}
\subsubsection*{\underline{Attributes}}
\texttt{nuclearSpinRef}: a label identifying the nuclear spin coupled to $\boldsymbol{J}$ (or $\boldsymbol{F_1}$) to form the total angular momentum.
\subsubsection*{\underline{Description}}
$F$ is the quantum number associated with the total angular momentum including nuclear spin: $\boldsymbol{F} = \boldsymbol{J} + \boldsymbol{I_1}$ if only one such coupling is resolved; $\boldsymbol{F} = \boldsymbol{F_1} + \boldsymbol{I_2}$ if both couplings are resolved.
\subsubsection*{\underline{Restrictions}}
\begin{itemize}
\item non-negative integer or half-integer.
\end{itemize}
}
}
\vspace{0.5cm}
\fbox{
\parbox{12cm}{
\subsection*{$r$}
\subsubsection*{\underline{Element}}
\texttt{<hunda:r>}
\subsubsection*{\underline{Attributes}}
\texttt{name}: a string identifying this ranking index.
\subsubsection*{\underline{Description}}
$r$ is a named, positive integer label identifying the state if no other good quantum numbers or symmetries are known.
\subsubsection*{\underline{Restrictions}}
\begin{itemize}
\item positive integer.
\end{itemize}
}
}
\vspace{0.5cm}
\fbox{
\parbox{12cm}{
\subsection*{\texttt{parity}}
\subsubsection*{\underline{Element}}
\texttt{<hunda:parity>}
\subsubsection*{\underline{Attributes}}
None.
\subsubsection*{\underline{Description}}
\texttt{parity} is the total parity: the parity of the total molecular wavefunction (excluding nuclear spin) with respect to inversion through the molecular centre of mass of all particles' coordinates in the laboratory coordinate system, the $\hat{E}^*$ operation.
\subsubsection*{\underline{Restrictions}}
\begin{itemize}
\item `$+$' or `$-$'.
\end{itemize}
}
}
\vspace{0.5cm}
\fbox{
\parbox{12cm}{
\subsection*{\texttt{kronigParity}}
\subsubsection*{\underline{Element}}
\texttt{<hunda:kronigParity>}
\subsubsection*{\underline{Attributes}}
None.
\subsubsection*{\underline{Description}}
\texttt{kronigParity} is the `rotationless' parity: the parity of the total molecular wavefunction excluding nuclear spin and rotation with respect to inversion through the molecular centre of mass of all particles' coordinates in the laboratory coordinate system. For integral $J$, the levels are called: \begin{align*}e\;\mathrm{if\;parity\;is}\;+(-1)^J,\\f\;\mathrm{if\;parity\;is}\;-(-1)^J.\end{align*}For half-odd integer $J$, the levels are called: \begin{align*}e\;\mathrm{if\;parity\;is}\;+(-1)^{J-\frac{1}{2}},\\f\;\mathrm{if\;parity\;is}\;-(-1)^{J-\frac{1}{2}}.\end{align*}.
\subsubsection*{\underline{Restrictions}}
\begin{itemize}
\item `e' or `f'.
\end{itemize}
}
}
\vspace{0.5cm}
\fbox{
\parbox{12cm}{
\subsection*{\texttt{asSym}}
\subsubsection*{\underline{Element}}
\texttt{<hunda:asSym>}
\subsubsection*{\underline{Attributes}}
None.
\subsubsection*{\underline{Description}}
\texttt{asSym} is (for diatomic molecules with a centre of inversion) the symmetry of the rovibronic wavefunction: `a' or `s' such that the total wavefunction including nuclear spin is symmetric or antisymmetric with respect to permutation of the identical nuclei ($\hat{P}_{12}$), according to whether they are bosons or fermions respectively.
\subsubsection*{\underline{Restrictions}}
\begin{itemize}
\item `s' or `a'.
\end{itemize}
}
}
\vspace{0.5cm}




\section*{Hund's case (b) diatomics: hundb}
\fbox{
\parbox{12cm}{
\subsection*{\texttt{ElecStateLabel}}
\subsubsection*{\underline{Element}}
\texttt{<hundb:ElecStateLabel>}
\subsubsection*{\underline{Attributes}}
None.
\subsubsection*{\underline{Description}}
\texttt{ElecStateLabel} is a label identifying the electronic state: $X$, $A$, $a$, $B$, etc..
\subsubsection*{\underline{Restrictions}}
\begin{itemize}
\item string.
\end{itemize}
}
}
\vspace{0.5cm}
\fbox{
\parbox{12cm}{
\subsection*{\texttt{elecInv}}
\subsubsection*{\underline{Element}}
\texttt{<hundb:elecInv>}
\subsubsection*{\underline{Attributes}}
None.
\subsubsection*{\underline{Description}}
\texttt{elecInv} is the parity of the electronic wavefunction with respect to inversion through the molecular centre of mass in the molecular coordinate system.
\subsubsection*{\underline{Restrictions}}
\begin{itemize}
\item `g' or `u'.
\end{itemize}
}
}
\vspace{0.5cm}
\fbox{
\parbox{12cm}{
\subsection*{\texttt{elecRefl}}
\subsubsection*{\underline{Element}}
\texttt{<hundb:elecRefl>}
\subsubsection*{\underline{Attributes}}
None.
\subsubsection*{\underline{Description}}
\texttt{elecRefl} is the parity of the electronic wavefunction with respect to reflection in a plane containing the molecular symmetry axis in the molecular coordinate system (equivalent to inversion through the molecular centre of mass in the laboratory coordinate system).
\subsubsection*{\underline{Restrictions}}
\begin{itemize}
\item `$+$' or `$-$'.
\end{itemize}
}
}
\vspace{0.5cm}
\fbox{
\parbox{12cm}{
\subsection*{$|\mathit{\Lambda}|$}
\subsubsection*{\underline{Element}}
\texttt{<hundb:Lambda>}
\subsubsection*{\underline{Attributes}}
None.
\subsubsection*{\underline{Description}}
$|\mathit{\Lambda}|$ is the quantum number associated with the magnitude of the projection of the total electronic orbital angular momentum, $\boldsymbol{L}$, onto the molecular axis.
\subsubsection*{\underline{Restrictions}}
\begin{itemize}
\item non-negative integer.
\end{itemize}
}
}
\vspace{0.5cm}
\fbox{
\parbox{12cm}{
\subsection*{$S$}
\subsubsection*{\underline{Element}}
\texttt{<hundb:S>}
\subsubsection*{\underline{Attributes}}
None.
\subsubsection*{\underline{Description}}
$S$ is the quantum number associated with the total electronic spin angular momentum, $\boldsymbol{S}$.
\subsubsection*{\underline{Restrictions}}
\begin{itemize}
\item non-negative integer or half-odd integer.
\end{itemize}
}
}
\vspace{0.5cm}
\fbox{
\parbox{12cm}{
\subsection*{$v$}
\subsubsection*{\underline{Element}}
\texttt{<hundb:v>}
\subsubsection*{\underline{Attributes}}
None.
\subsubsection*{\underline{Description}}
$v$ is the vibrational quantum number.
\subsubsection*{\underline{Restrictions}}
\begin{itemize}
\item non-negative integer.
\end{itemize}
}
}
\vspace{0.5cm}
\fbox{
\parbox{12cm}{
\subsection*{$J$}
\subsubsection*{\underline{Element}}
\texttt{<hundb:J>}
\subsubsection*{\underline{Attributes}}
None.
\subsubsection*{\underline{Description}}
$J$ is the quantum number associated with the total angular momentum excluding nuclear spin: $\boldsymbol{J} = \boldsymbol{N} + \boldsymbol{S} = \boldsymbol{L} + \boldsymbol{S} + \boldsymbol{R}$.
\subsubsection*{\underline{Restrictions}}
\begin{itemize}
\item non-negative integer or half-integer.
\item $|N - S| \le J \le N + S$.
\end{itemize}
}
}
\vspace{0.5cm}
\fbox{
\parbox{12cm}{
\subsection*{$N$}
\subsubsection*{\underline{Element}}
\texttt{<hundb:N>}
\subsubsection*{\underline{Attributes}}
None.
\subsubsection*{\underline{Description}}
$N$ is the quantum number associated with the total angular momentum excluding electronic and nuclear spin, $\boldsymbol{N}$: $\boldsymbol{J} = \boldsymbol{N} + \boldsymbol{S}$.
\subsubsection*{\underline{Restrictions}}
\begin{itemize}
\item non-negative integer.
\item $N \ge |\mathit{\Lambda}|$.
\end{itemize}
}
}
\vspace{0.5cm}
\fbox{
\parbox{12cm}{
\subsection*{\texttt{SpinComponentLabel}}
\subsubsection*{\underline{Element}}
\texttt{<hundb:SpinComponentLabel>}
\subsubsection*{\underline{Attributes}}
None.
\subsubsection*{\underline{Description}}
\texttt{SpinComponentLabel} is the positive integer identifying the spin-component label, $F_x$, where $x = 1,2,3, \cdots$ in order of increasing energy for a given value of $J$ - see Herzberg, \emph{Spectra of Diatomic Molecules}, Van Nostrand, Princeton, N.J., 1950.
\subsubsection*{\underline{Restrictions}}
\begin{itemize}
\item positive integer.
\end{itemize}
}
}
\vspace{0.5cm}
\fbox{
\parbox{12cm}{
\subsection*{$F_1$}
\subsubsection*{\underline{Element}}
\texttt{<hundb:F1>}
\subsubsection*{\underline{Attributes}}
\texttt{nuclearSpinRef}: a label identifying the nuclear spin coupled to $\boldsymbol{J}$ to form the intermediate angular momentum.
\subsubsection*{\underline{Description}}
$F_1$ is the intermediate angular momentum quantum number associated with the coupling of the rotational angular momentum and nuclear spin of nucleus 1 where two such couplings are resolved: $\boldsymbol{F_1} = \boldsymbol{J} + \boldsymbol{I_1}$; $F_1$ is often not a good quantum number.
\subsubsection*{\underline{Restrictions}}
\begin{itemize}
\item non-negative integer or half-integer.
\item $|J - I_1| \le F_1 \le J + I_1$.
\end{itemize}
}
}
\vspace{0.5cm}
\fbox{
\parbox{12cm}{
\subsection*{$F$}
\subsubsection*{\underline{Element}}
\texttt{<hundb:F>}
\subsubsection*{\underline{Attributes}}
\texttt{nuclearSpinRef}: a label identifying the nuclear spin coupled to $\boldsymbol{J}$ (or $\boldsymbol{F_1}$) to form the total angular momentum.
\subsubsection*{\underline{Description}}
$F$ is the quantum number associated with the total angular momentum including nuclear spin: $\boldsymbol{F} = \boldsymbol{J} + \boldsymbol{I_1}$ if only one such coupling is resolved; $\boldsymbol{F} = \boldsymbol{F_1} + \boldsymbol{I_2}$ if both couplings are resolved.
\subsubsection*{\underline{Restrictions}}
\begin{itemize}
\item non-negative integer or half-integer.
\end{itemize}
}
}
\vspace{0.5cm}
\fbox{
\parbox{12cm}{
\subsection*{$r$}
\subsubsection*{\underline{Element}}
\texttt{<hundb:r>}
\subsubsection*{\underline{Attributes}}
\texttt{name}: a string identifying this ranking index.
\subsubsection*{\underline{Description}}
$r$ is a named, positive integer label identifying the state if no other good quantum numbers or symmetries are known.
\subsubsection*{\underline{Restrictions}}
\begin{itemize}
\item positive integer.
\end{itemize}
}
}
\vspace{0.5cm}
\fbox{
\parbox{12cm}{
\subsection*{\texttt{parity}}
\subsubsection*{\underline{Element}}
\texttt{<hundb:parity>}
\subsubsection*{\underline{Attributes}}
None.
\subsubsection*{\underline{Description}}
\texttt{parity} is the total parity: the parity of the total molecular wavefunction (excluding nuclear spin) with respect to inversion through the molecular centre of mass of all particles' coordinates in the laboratory coordinate system, the $\hat{E}^*$ operation.
\subsubsection*{\underline{Restrictions}}
\begin{itemize}
\item `$+$' or `$-$'.
\end{itemize}
}
}
\vspace{0.5cm}
\fbox{
\parbox{12cm}{
\subsection*{\texttt{kronigParity}}
\subsubsection*{\underline{Element}}
\texttt{<hundb:kronigParity>}
\subsubsection*{\underline{Attributes}}
None.
\subsubsection*{\underline{Description}}
\texttt{kronigParity} is the `rotationless' parity: the parity of the total molecular wavefunction excluding nuclear spin and rotation with respect to inversion through the molecular centre of mass of all particles' coordinates in the laboratory coordinate system. For integral $J$, the levels are called: \begin{align*}e\;\mathrm{if\;parity\;is}\;+(-1)^J,\\f\;\mathrm{if\;parity\;is}\;-(-1)^J.\end{align*}For half-odd integer $J$, the levels are called: \begin{align*}e\;\mathrm{if\;parity\;is}\;+(-1)^{J-\frac{1}{2}},\\f\;\mathrm{if\;parity\;is}\;-(-1)^{J-\frac{1}{2}}.\end{align*}.
\subsubsection*{\underline{Restrictions}}
\begin{itemize}
\item `e' or `f'.
\end{itemize}
}
}
\vspace{0.5cm}
\fbox{
\parbox{12cm}{
\subsection*{\texttt{asSym}}
\subsubsection*{\underline{Element}}
\texttt{<hundb:asSym>}
\subsubsection*{\underline{Attributes}}
None.
\subsubsection*{\underline{Description}}
\texttt{asSym} is (for diatomic molecules with a centre of inversion) the symmetry of the rovibronic wavefunction: `a' or `s' such that the total wavefunction including nuclear spin is symmetric or antisymmetric with respect to permutation of the identical nuclei ($\hat{P}_{12}$), according to whether they are bosons or fermions respectively.
\subsubsection*{\underline{Restrictions}}
\begin{itemize}
\item `s' or `a'.
\end{itemize}
}
}
\vspace{0.5cm}




\section*{Closed-shell, linear triatomic molecules: ltcs}
\fbox{
\parbox{12cm}{
\subsection*{\texttt{ElecStateLabel}}
\subsubsection*{\underline{Element}}
\texttt{<ltcs:ElecStateLabel>}
\subsubsection*{\underline{Attributes}}
None.
\subsubsection*{\underline{Description}}
\texttt{ElecStateLabel} is a label identifying the electronic state: $X$, $A$, $a$, $B$, etc..
\subsubsection*{\underline{Restrictions}}
\begin{itemize}
\item string.
\end{itemize}
}
}
\vspace{0.5cm}
\fbox{
\parbox{12cm}{
\subsection*{$v_1$}
\subsubsection*{\underline{Element}}
\texttt{<ltcs:v1>}
\subsubsection*{\underline{Attributes}}
None.
\subsubsection*{\underline{Description}}
$v_1$ is the vibrational quantum number associated with the $\nu_1$ normal mode.
\subsubsection*{\underline{Restrictions}}
\begin{itemize}
\item non-negative integer.
\end{itemize}
}
}
\vspace{0.5cm}
\fbox{
\parbox{12cm}{
\subsection*{$v_2$}
\subsubsection*{\underline{Element}}
\texttt{<ltcs:v2>}
\subsubsection*{\underline{Attributes}}
None.
\subsubsection*{\underline{Description}}
$v_2$ is the vibrational quantum number associated with the doubly-degenerate $\nu_2$ normal mode.
\subsubsection*{\underline{Restrictions}}
\begin{itemize}
\item non-negative integer.
\end{itemize}
}
}
\vspace{0.5cm}
\fbox{
\parbox{12cm}{
\subsection*{$l_2$}
\subsubsection*{\underline{Element}}
\texttt{<ltcs:l2>}
\subsubsection*{\underline{Attributes}}
None.
\subsubsection*{\underline{Description}}
$l_2$ is the vibrational angular momentum quantum number associated with the degenerate bending vibration, $\nu_2$; positive and negative values distinguish $l$-type doubling components.
\subsubsection*{\underline{Restrictions}}
\begin{itemize}
\item integer.
\item $|l_2| = v_2, v_2 - 2, \cdots, 1 \;\mathrm{or}\;0$.
\end{itemize}
}
}
\vspace{0.5cm}
\fbox{
\parbox{12cm}{
\subsection*{$v_3$}
\subsubsection*{\underline{Element}}
\texttt{<ltcs:v3>}
\subsubsection*{\underline{Attributes}}
None.
\subsubsection*{\underline{Description}}
$v_3$ is the vibrational quantum number associated with the $\nu_3$ normal mode.
\subsubsection*{\underline{Restrictions}}
\begin{itemize}
\item non-negative integer.
\end{itemize}
}
}
\vspace{0.5cm}
\fbox{
\parbox{12cm}{
\subsection*{$J$}
\subsubsection*{\underline{Element}}
\texttt{<ltcs:J>}
\subsubsection*{\underline{Attributes}}
None.
\subsubsection*{\underline{Description}}
$J$ is the quantum number associated with the total angular momentum excluding nuclear spin, $\boldsymbol{J}$.
\subsubsection*{\underline{Restrictions}}
\begin{itemize}
\item non-negative integer.
\end{itemize}
}
}
\vspace{0.5cm}
\fbox{
\parbox{12cm}{
\subsection*{$F_1$}
\subsubsection*{\underline{Element}}
\texttt{<ltcs:F1>}
\subsubsection*{\underline{Attributes}}
\texttt{nuclearSpinRef}: a label identifying the nuclear spin coupled to $\boldsymbol{J}$ to form the intermediate angular momentum.
\subsubsection*{\underline{Description}}
$F_1$ is the intermediate angular momentum quantum number associated with the coupling of the rotational angular momentum and nuclear spin of nucleus 1 where two or more such couplings are resolved: $\boldsymbol{F_1} = \boldsymbol{J} + \boldsymbol{I_1}$; $F_1$ is often not a good quantum number.
\subsubsection*{\underline{Restrictions}}
\begin{itemize}
\item non-negative integer or half-integer.
\item $|J - I_1| \le F_1 \le J + I_1$.
\end{itemize}
}
}
\vspace{0.5cm}
\fbox{
\parbox{12cm}{
\subsection*{$F_2$}
\subsubsection*{\underline{Element}}
\texttt{<ltcs:F2>}
\subsubsection*{\underline{Attributes}}
\texttt{nuclearSpinRef}: a label identifying the nuclear spin coupled to $\boldsymbol{F_1}$ to form an intermediate angular momentum.
\subsubsection*{\underline{Description}}
$F_2$ is the intermediate angular momentum quantum number associated with the coupling of the rotational angular momentum and nuclear spin of nucleus 2 where three such couplings are resolved: $\boldsymbol{F_2} = \boldsymbol{F_1} + \boldsymbol{I_2}$; $F_2$ is often not a good quantum number.
\subsubsection*{\underline{Restrictions}}
\begin{itemize}
\item non-negative integer or half-integer.
\item $|F_1 - I_2| \le F_2 \le F_1 + I_2$.
\end{itemize}
}
}
\vspace{0.5cm}
\fbox{
\parbox{12cm}{
\subsection*{$F$}
\subsubsection*{\underline{Element}}
\texttt{<ltcs:F>}
\subsubsection*{\underline{Attributes}}
\texttt{nuclearSpinRef}: a label identifying the nuclear spin coupled to $\boldsymbol{J}$, $\boldsymbol{F_1}$, or $\boldsymbol{F}$ to form the total angular momentum.
\subsubsection*{\underline{Description}}
$F$ is the quantum number associated with the total angular momentum including nuclear spin: $\boldsymbol{F} = \boldsymbol{J} + \boldsymbol{I_1}$ if only one hyperfine coupling is resolved, $\boldsymbol{F} = \boldsymbol{F_1} + \boldsymbol{I_2}$ if two couplings are resolved, or $\boldsymbol{F} = \boldsymbol{F_2} + \boldsymbol{I_3}$ if all three couplings are resolved.
\subsubsection*{\underline{Restrictions}}
\begin{itemize}
\item non-negative integer or half-integer.
\item $|F_2 - I_3| \le F \le F_2 + I_3$.
\end{itemize}
}
}
\vspace{0.5cm}
\fbox{
\parbox{12cm}{
\subsection*{$r$}
\subsubsection*{\underline{Element}}
\texttt{<ltcs:r>}
\subsubsection*{\underline{Attributes}}
\texttt{name}: a string identifying this ranking index.
\subsubsection*{\underline{Description}}
$r$ is a named, positive integer label identifying the state if no other good quantum numbers or symmetries are known.
\subsubsection*{\underline{Restrictions}}
\begin{itemize}
\item positive integer.
\end{itemize}
}
}
\vspace{0.5cm}
\fbox{
\parbox{12cm}{
\subsection*{\texttt{parity}}
\subsubsection*{\underline{Element}}
\texttt{<ltcs:parity>}
\subsubsection*{\underline{Attributes}}
None.
\subsubsection*{\underline{Description}}
\texttt{parity} is the total parity: the parity of the total molecular wavefunction (excluding nuclear spin) with respect to inversion through the molecular centre of mass of all particles' coordinates in the laboratory coordinate system, the $\hat{E}^*$ operation.
\subsubsection*{\underline{Restrictions}}
\begin{itemize}
\item `$+$' or `$-$'.
\end{itemize}
}
}
\vspace{0.5cm}
\fbox{
\parbox{12cm}{
\subsection*{\texttt{kronigParity}}
\subsubsection*{\underline{Element}}
\texttt{<ltcs:kronigParity>}
\subsubsection*{\underline{Attributes}}
None.
\subsubsection*{\underline{Description}}
\texttt{kronigParity} is the `rotationless' parity: the parity of the total molecular wavefunction excluding nuclear spin and rotation with respect to inversion through the molecular centre of mass of all particles' coordinates in the laboratory coordinate system. For integral $J$, the levels are called: \begin{align*}e\;\mathrm{if\;parity\;is}\;+(-1)^J,\\f\;\mathrm{if\;parity\;is}\;-(-1)^J.\end{align*}For half-odd integer $J$, the levels are called: \begin{align*}e\;\mathrm{if\;parity\;is}\;+(-1)^{J-\frac{1}{2}},\\f\;\mathrm{if\;parity\;is}\;-(-1)^{J-\frac{1}{2}}.\end{align*}.
\subsubsection*{\underline{Restrictions}}
\begin{itemize}
\item `e' or `f'.
\end{itemize}
}
}
\vspace{0.5cm}
\fbox{
\parbox{12cm}{
\subsection*{\texttt{asSym}}
\subsubsection*{\underline{Element}}
\texttt{<ltcs:asSym>}
\subsubsection*{\underline{Attributes}}
None.
\subsubsection*{\underline{Description}}
\texttt{asSym} is (for linear molecules with a centre of inversion) the symmetry of the rovibronic wavefunction: `a' or `s' such that the total wavefunction including nuclear spin is symmetric or antisymmetric with respect to permutation of the identical nuclei ($\hat{P}_{12}$), according to whether they are bosons or fermions respectively.
\subsubsection*{\underline{Restrictions}}
\begin{itemize}
\item `s' or `a'.
\end{itemize}
}
}
\vspace{0.5cm}




\section*{Closed-shell, non-linear triatomics: nltcs}
\fbox{
\parbox{12cm}{
\subsection*{\texttt{ElecStateLabel}}
\subsubsection*{\underline{Element}}
\texttt{<nltcs:ElecStateLabel>}
\subsubsection*{\underline{Attributes}}
None.
\subsubsection*{\underline{Description}}
\texttt{ElecStateLabel} is a label identifying the electronic state: $X$, $A$, $a$, $B$, etc..
\subsubsection*{\underline{Restrictions}}
\begin{itemize}
\item string.
\end{itemize}
}
}
\vspace{0.5cm}
\fbox{
\parbox{12cm}{
\subsection*{$v_1$}
\subsubsection*{\underline{Element}}
\texttt{<nltcs:v1>}
\subsubsection*{\underline{Attributes}}
None.
\subsubsection*{\underline{Description}}
$v_1$ is the vibrational quantum number associated with the $\nu_1$ normal mode.
\subsubsection*{\underline{Restrictions}}
\begin{itemize}
\item non-negative integer.
\end{itemize}
}
}
\vspace{0.5cm}
\fbox{
\parbox{12cm}{
\subsection*{$v_2$}
\subsubsection*{\underline{Element}}
\texttt{<nltcs:v2>}
\subsubsection*{\underline{Attributes}}
None.
\subsubsection*{\underline{Description}}
$v_2$ is the vibrational quantum number associated with the $\nu_2$ normal mode.
\subsubsection*{\underline{Restrictions}}
\begin{itemize}
\item non-negative integer.
\end{itemize}
}
}
\vspace{0.5cm}
\fbox{
\parbox{12cm}{
\subsection*{$v_3$}
\subsubsection*{\underline{Element}}
\texttt{<nltcs:v3>}
\subsubsection*{\underline{Attributes}}
None.
\subsubsection*{\underline{Description}}
$v_3$ is the vibrational quantum number associated with the $\nu_3$ normal mode.
\subsubsection*{\underline{Restrictions}}
\begin{itemize}
\item non-negative integer.
\end{itemize}
}
}
\vspace{0.5cm}
\fbox{
\parbox{12cm}{
\subsection*{$J$}
\subsubsection*{\underline{Element}}
\texttt{<nltcs:J>}
\subsubsection*{\underline{Attributes}}
None.
\subsubsection*{\underline{Description}}
$J$ is the quantum number associated with the total angular momentum excluding nuclear spin, $\boldsymbol{J}$.
\subsubsection*{\underline{Restrictions}}
\begin{itemize}
\item non-negative integer.
\end{itemize}
}
}
\vspace{0.5cm}
\fbox{
\parbox{12cm}{
\subsection*{$K_a$}
\subsubsection*{\underline{Element}}
\texttt{<nltcs:Ka>}
\subsubsection*{\underline{Attributes}}
None.
\subsubsection*{\underline{Description}}
$K_a$ is the rotational quantum label of an asymmetric top molecule, correlating to $K$ in the prolate symmetric top limit.
\subsubsection*{\underline{Restrictions}}
\begin{itemize}
\item non-negative integer.
\item $K_a \le J$.
\item $K_a + K_c = J \;\mathrm{or}\;J + 1$.
\end{itemize}
}
}
\vspace{0.5cm}
\fbox{
\parbox{12cm}{
\subsection*{$K_c$}
\subsubsection*{\underline{Element}}
\texttt{<nltcs:Kc>}
\subsubsection*{\underline{Attributes}}
None.
\subsubsection*{\underline{Description}}
$K_c$ is the rotational quantum label of an asymmetric top molecule, correlating to $K$ in the oblate symmetric top limit.
\subsubsection*{\underline{Restrictions}}
\begin{itemize}
\item non-negative integer.
\item $K_c \le J$.
\item $K_a + K_c = J \;\mathrm{or}\;J + 1$.
\end{itemize}
}
}
\vspace{0.5cm}
\fbox{
\parbox{12cm}{
\subsection*{$F_1$}
\subsubsection*{\underline{Element}}
\texttt{<nltcs:F1>}
\subsubsection*{\underline{Attributes}}
\texttt{nuclearSpinRef}: a label identifying the nuclear spin coupled to $\boldsymbol{J}$ to form the intermediate angular momentum.
\subsubsection*{\underline{Description}}
$F_1$ is the intermediate angular momentum quantum number associated with the coupling of the rotational angular momentum and nuclear spin of nucleus 1 where two or more such couplings are resolved: $\boldsymbol{F_1} = \boldsymbol{J} + \boldsymbol{I_1}$; $F_1$ is often not a good quantum number.
\subsubsection*{\underline{Restrictions}}
\begin{itemize}
\item non-negative integer or half-integer.
\item $|J - I_1| \le F_1 \le J + I_1$.
\end{itemize}
}
}
\vspace{0.5cm}
\fbox{
\parbox{12cm}{
\subsection*{$F_2$}
\subsubsection*{\underline{Element}}
\texttt{<nltcs:F2>}
\subsubsection*{\underline{Attributes}}
\texttt{nuclearSpinRef}: a label identifying the nuclear spin coupled to $\boldsymbol{F_1}$ to form an intermediate angular momentum.
\subsubsection*{\underline{Description}}
$F_2$ is the intermediate angular momentum quantum number associated with the coupling of the rotational angular momentum and nuclear spin of nucleus 2 where three such couplings are resolved: $\boldsymbol{F_2} = \boldsymbol{F_1} + \boldsymbol{I_2}$; $F_2$ is often not a good quantum number.
\subsubsection*{\underline{Restrictions}}
\begin{itemize}
\item non-negative integer or half-integer.
\item $|F_1 - I_2| \le F_2 \le F_1 + I_2$.
\end{itemize}
}
}
\vspace{0.5cm}
\fbox{
\parbox{12cm}{
\subsection*{$F$}
\subsubsection*{\underline{Element}}
\texttt{<nltcs:F>}
\subsubsection*{\underline{Attributes}}
\texttt{nuclearSpinRef}: a label identifying the nuclear spin coupled to $\boldsymbol{J}$, $\boldsymbol{F_1}$, or $\boldsymbol{F}$ to form the total angular momentum.
\subsubsection*{\underline{Description}}
$F$ is the quantum number associated with the total angular momentum including nuclear spin: $\boldsymbol{F} = \boldsymbol{J} + \boldsymbol{I_1}$ if only one hyperfine coupling is resolved, $\boldsymbol{F} = \boldsymbol{F_1} + \boldsymbol{I_2}$ if two couplings are resolved, or $\boldsymbol{F} = \boldsymbol{F_2} + \boldsymbol{I_3}$ if all three couplings are resolved.
\subsubsection*{\underline{Restrictions}}
\begin{itemize}
\item non-negative integer or half-integer.
\item $|F_2 - I_3| \le F \le F_2 + I_3$.
\end{itemize}
}
}
\vspace{0.5cm}
\fbox{
\parbox{12cm}{
\subsection*{$r$}
\subsubsection*{\underline{Element}}
\texttt{<nltcs:r>}
\subsubsection*{\underline{Attributes}}
\texttt{name}: a string identifying this ranking index.
\subsubsection*{\underline{Description}}
$r$ is a named, positive integer label identifying the state if no other good quantum numbers or symmetries are known.
\subsubsection*{\underline{Restrictions}}
\begin{itemize}
\item positive integer.
\end{itemize}
}
}
\vspace{0.5cm}
\fbox{
\parbox{12cm}{
\subsection*{\texttt{parity}}
\subsubsection*{\underline{Element}}
\texttt{<nltcs:parity>}
\subsubsection*{\underline{Attributes}}
None.
\subsubsection*{\underline{Description}}
\texttt{parity} is the total parity: the parity of the total molecular wavefunction (excluding nuclear spin) with respect to inversion through the molecular centre of mass of all particles' coordinates in the laboratory coordinate system, the $\hat{E}^*$ operation.
\subsubsection*{\underline{Restrictions}}
\begin{itemize}
\item `$+$' or `$-$'.
\end{itemize}
}
}
\vspace{0.5cm}
\fbox{
\parbox{12cm}{
\subsection*{\texttt{asSym}}
\subsubsection*{\underline{Element}}
\texttt{<nltcs:asSym>}
\subsubsection*{\underline{Attributes}}
None.
\subsubsection*{\underline{Description}}
\texttt{asSym} is the symmetry of the rovibronic wavefunction: `a' or `s' such that the total wavefunction including nuclear spin is symmetric or antisymmetric with respect to permutation of the identical nuclei ($\hat{P}_{12}$), according to whether they are bosons or fermions respectively.
\subsubsection*{\underline{Restrictions}}
\begin{itemize}
\item `s' or `a'.
\end{itemize}
}
}
\vspace{0.5cm}




\section*{Closed shell, symmetric-top molecules: stcs}
\fbox{
\parbox{12cm}{
\subsection*{\texttt{ElecStateLabel}}
\subsubsection*{\underline{Element}}
\texttt{<stcs:ElecStateLabel>}
\subsubsection*{\underline{Attributes}}
None.
\subsubsection*{\underline{Description}}
\texttt{ElecStateLabel} is a label identifying the electronic state.
\subsubsection*{\underline{Restrictions}}
\begin{itemize}
\item string.
\end{itemize}
}
}
\vspace{0.5cm}
\fbox{
\parbox{12cm}{
\subsection*{$v_i$}
\subsubsection*{\underline{Element}}
\texttt{<stcs:vi>}
\subsubsection*{\underline{Attributes}}
\texttt{mode}: a positive integer, identifying the normal mode that this quantum number is associated with.
\subsubsection*{\underline{Description}}
$v_i$ is the vibrational quantum number associated with the $\nu_i$ normal mode.
\subsubsection*{\underline{Restrictions}}
\begin{itemize}
\item non-negative integer.
\end{itemize}
}
}
\vspace{0.5cm}
\fbox{
\parbox{12cm}{
\subsection*{$l_i$}
\subsubsection*{\underline{Element}}
\texttt{<stcs:li>}
\subsubsection*{\underline{Attributes}}
\texttt{mode}: a positive integer, identifying the degenerate normal mode that this vibrational angular momentum quantum number is associated with.
\subsubsection*{\underline{Description}}
$l_i$ is the vibrational angular momentum quantum number associated with the degenerate $\nu_i$ normal mode; positive and negative values distinguish $l$-type doubling components.
\subsubsection*{\underline{Restrictions}}
\begin{itemize}
\item non-negative integer.
\item $|l_i| = v_i, v_i-2, \cdots, 1\;\mathrm{or}\;0$.
\end{itemize}
}
}
\vspace{0.5cm}
\fbox{
\parbox{12cm}{
\subsection*{\texttt{vibInv}}
\subsubsection*{\underline{Element}}
\texttt{<stcs:vibInv>}
\subsubsection*{\underline{Attributes}}
None.
\subsubsection*{\underline{Description}}
\texttt{vibInv} is the parity of the vibrational wavefunction with respect to inversion through the molecular centre of mass in the molecular coordinate system. Only really necessary for molecules with a low barrier to such an inversion (for example, $\mathrm{NH_3}$).
\subsubsection*{\underline{Restrictions}}
\begin{itemize}
\item `s' or `a'.
\end{itemize}
}
}
\vspace{0.5cm}
\fbox{
\parbox{12cm}{
\subsection*{\texttt{vibSym}}
\subsubsection*{\underline{Element}}
\texttt{<stcs:vibSym>}
\subsubsection*{\underline{Attributes}}
\texttt{group}: the symmetry group to which this species belongs.
\subsubsection*{\underline{Description}}
\texttt{vibSym} is the symmetry species of the vibrational wavefunction, in some appropriate symmetry group.
\subsubsection*{\underline{Restrictions}}
\begin{itemize}
\item string.
\end{itemize}
}
}
\vspace{0.5cm}
\fbox{
\parbox{12cm}{
\subsection*{$J$}
\subsubsection*{\underline{Element}}
\texttt{<stcs:J>}
\subsubsection*{\underline{Attributes}}
None.
\subsubsection*{\underline{Description}}
$J$ is the quantum number associated with the total angular momentum excluding nuclear spin, $\boldsymbol{J}$.
\subsubsection*{\underline{Restrictions}}
\begin{itemize}
\item non-negative integer.
\end{itemize}
}
}
\vspace{0.5cm}
\fbox{
\parbox{12cm}{
\subsection*{$K$}
\subsubsection*{\underline{Element}}
\texttt{<stcs:K>}
\subsubsection*{\underline{Attributes}}
None.
\subsubsection*{\underline{Description}}
$K$ is the quantum number associated with the projection of the total angular momentum excluding nuclear spin, $\boldsymbol{J}$, onto the molecular symmetry axis.
\subsubsection*{\underline{Restrictions}}
\begin{itemize}
\item non-negative integer.
\item $K \le J$.
\end{itemize}
}
}
\vspace{0.5cm}
\fbox{
\parbox{12cm}{
\subsection*{\texttt{rotSym}}
\subsubsection*{\underline{Element}}
\texttt{<stcs:rotSym>}
\subsubsection*{\underline{Attributes}}
\texttt{group}: the symmetry group to which this species belongs.
\subsubsection*{\underline{Description}}
\texttt{rotSym} is the symmetry species of the rotational wavefunction, in some appropriate symmetry group.
\subsubsection*{\underline{Restrictions}}
\begin{itemize}
\item string.
\end{itemize}
}
}
\vspace{0.5cm}
\fbox{
\parbox{12cm}{
\subsection*{$I$}
\subsubsection*{\underline{Element}}
\texttt{<stcs:I>}
\subsubsection*{\underline{Attributes}}
\texttt{nuclearSpinRef}: a label identifying the group of nuclear spins coupled to one another to form a total nuclear spin angular momentum.
\subsubsection*{\underline{Description}}
$I$ is the quantum number associated with the total nuclear spin angular momentum: $\boldsymbol{I} = \boldsymbol{I_1} + \boldsymbol{I_2} + \cdots$ where nuclei $1, 2, \cdots$ have individual nuclear spin angular momenta $\boldsymbol{I_1}, \boldsymbol{I_2}, \cdots$.
\subsubsection*{\underline{Restrictions}}
\begin{itemize}
\item non-negative integer or half-integer.
\end{itemize}
}
}
\vspace{0.5cm}
\fbox{
\parbox{12cm}{
\subsection*{$F_j$}
\subsubsection*{\underline{Element}}
\texttt{<stcs:Fj>}
\subsubsection*{\underline{Attributes}}
\begin{itemize}\item\texttt{nuclearSpinRef}: a label identifying the nuclear spin being coupled to $\boldsymbol{J}$ or $\boldsymbol{F_{j-1}}$ to form an intermediate angular momentum;\item $j$: an integer label identifying the order of the hyperfine coupling\end{itemize}.
\subsubsection*{\underline{Description}}
$F_j$ is the intermediate angular momentum quantum number associated with the coupling of the nuclear spin angular momentum of nucleus $j$ to the intermediate angular momentum: $\boldsymbol{F_1} = \boldsymbol{J} + \boldsymbol{I_1}$ or $\boldsymbol{F_j} = \boldsymbol{F_{j-1}} + \boldsymbol{I_j}$; $\boldsymbol{F_j}$ is often not a good quantum number.
\subsubsection*{\underline{Restrictions}}
\begin{itemize}
\item non-negative integer or half-integer.
\end{itemize}
}
}
\vspace{0.5cm}
\fbox{
\parbox{12cm}{
\subsection*{$F$}
\subsubsection*{\underline{Element}}
\texttt{<stcs:F>}
\subsubsection*{\underline{Attributes}}
\texttt{nuclearSpinRef}: a label identifying the nuclear spin being coupled to $\boldsymbol{J}$ or $\boldsymbol{F_j}$ to form the total angular momentum.
\subsubsection*{\underline{Description}}
$F$ is the quantum number associated with the total angular momentum including nuclear spin: $\boldsymbol{F} = \boldsymbol{J} + \boldsymbol{I_1}$ if only one such coupling is resolved, $\boldsymbol{F} = \boldsymbol{F_{j-1}} + \boldsymbol{I_j}$ if two or more such couplings are resolved.
\subsubsection*{\underline{Restrictions}}
\begin{itemize}
\item non-negative integer or half-integer.
\end{itemize}
}
}
\vspace{0.5cm}
\fbox{
\parbox{12cm}{
\subsection*{$r$}
\subsubsection*{\underline{Element}}
\texttt{<stcs:r>}
\subsubsection*{\underline{Attributes}}
\texttt{name}: a string identifying this ranking index.
\subsubsection*{\underline{Description}}
$r$ is a named, positive integer label identifying the state if no other good quantum numbers or symmetries are known.
\subsubsection*{\underline{Restrictions}}
\begin{itemize}
\item positive integer.
\end{itemize}
}
}
\vspace{0.5cm}
\fbox{
\parbox{12cm}{
\subsection*{\texttt{parity}}
\subsubsection*{\underline{Element}}
\texttt{<stcs:parity>}
\subsubsection*{\underline{Attributes}}
None.
\subsubsection*{\underline{Description}}
\texttt{parity} is the total parity: the parity of the total molecular wavefunction (excluding nuclear spin) with respect to inversion through the molecular centre of mass of all particles' coordinates in the laboratory coordinate system, the $\hat{E}^*$ operation.
\subsubsection*{\underline{Restrictions}}
\begin{itemize}
\item `$+$' or `$-$'.
\end{itemize}
}
}
\vspace{0.5cm}




\section*{Closed-shell, linear polyatomic molecules: lpcs}
\fbox{
\parbox{12cm}{
\subsection*{\texttt{ElecStateLabel}}
\subsubsection*{\underline{Element}}
\texttt{<lpcs:ElecStateLabel>}
\subsubsection*{\underline{Attributes}}
None.
\subsubsection*{\underline{Description}}
\texttt{ElecStateLabel} is a label identifying the electronic state.
\subsubsection*{\underline{Restrictions}}
\begin{itemize}
\item string.
\end{itemize}
}
}
\vspace{0.5cm}
\fbox{
\parbox{12cm}{
\subsection*{$v_i$}
\subsubsection*{\underline{Element}}
\texttt{<lpcs:vi>}
\subsubsection*{\underline{Attributes}}
\texttt{mode}: a positive integer, identifying the normal mode that this quantum number is associated with.
\subsubsection*{\underline{Description}}
$v_i$ is the vibrational quantum number associated with the $\nu_i$ normal mode.
\subsubsection*{\underline{Restrictions}}
\begin{itemize}
\item non-negative integer.
\end{itemize}
}
}
\vspace{0.5cm}
\fbox{
\parbox{12cm}{
\subsection*{$l_i$}
\subsubsection*{\underline{Element}}
\texttt{<lpcs:li>}
\subsubsection*{\underline{Attributes}}
\texttt{mode}: a positive integer, identifying the degenerate normal mode that this vibrational angular momentum quantum number is associated with.
\subsubsection*{\underline{Description}}
$l_i$ is the vibrational angular momentum quantum number associated with the degenerate $\nu_i$ normal mode; positive and negative values distinguish $l$-type doubling components; if two or more degenerate vibrations are excited, $l_i$ is only approximately defined (\emph{i.e.} it is not a totally good quantum number) - see \emph{e.g.} Herzerg II, p.212.
\subsubsection*{\underline{Restrictions}}
\begin{itemize}
\item non-negative integer.
\item $|l_i| = v_i, v_i-2, \cdots, 1 \;\mathrm{or}\;0$.
\end{itemize}
}
}
\vspace{0.5cm}
\fbox{
\parbox{12cm}{
\subsection*{$l$}
\subsubsection*{\underline{Element}}
\texttt{<lpcs:l>}
\subsubsection*{\underline{Attributes}}
None.
\subsubsection*{\underline{Description}}
$l$ is the total vibrational angular momentum quantum number associated with resultant vibrational angular momentum about the internuclear axis.
\subsubsection*{\underline{Restrictions}}
\begin{itemize}
\item non-negative integer.
\end{itemize}
}
}
\vspace{0.5cm}
\fbox{
\parbox{12cm}{
\subsection*{\texttt{vibInv}}
\subsubsection*{\underline{Element}}
\texttt{<lpcs:vibInv>}
\subsubsection*{\underline{Attributes}}
None.
\subsubsection*{\underline{Description}}
\texttt{vibInv} is the parity of the vibrational wavefunction with respect to inversion through the molecular centre of mass in the molecular coordinate system.
\subsubsection*{\underline{Restrictions}}
\begin{itemize}
\item `g' or `u'.
\end{itemize}
}
}
\vspace{0.5cm}
\fbox{
\parbox{12cm}{
\subsection*{\texttt{vibRefl}}
\subsubsection*{\underline{Element}}
\texttt{<lpcs:vibRefl>}
\subsubsection*{\underline{Attributes}}
None.
\subsubsection*{\underline{Description}}
\texttt{vibRefl} is the parity of the vibrational wavefunction with respect to reflection in a plane containing the molecular symmetry axis in the molecular coordinate system.
\subsubsection*{\underline{Restrictions}}
\begin{itemize}
\item `$+$' or `$-$'.
\end{itemize}
}
}
\vspace{0.5cm}
\fbox{
\parbox{12cm}{
\subsection*{$J$}
\subsubsection*{\underline{Element}}
\texttt{<lpcs:J>}
\subsubsection*{\underline{Attributes}}
None.
\subsubsection*{\underline{Description}}
$J$ is the quantum number associated with the total angular momentum excluding nuclear spin, $\boldsymbol{J}$.
\subsubsection*{\underline{Restrictions}}
\begin{itemize}
\item non-negative integer.
\item $J \ge |l|$.
\end{itemize}
}
}
\vspace{0.5cm}
\fbox{
\parbox{12cm}{
\subsection*{$I$}
\subsubsection*{\underline{Element}}
\texttt{<lpcs:I>}
\subsubsection*{\underline{Attributes}}
\texttt{nuclearSpinRef}: a label, matching \texttt{/Q.+/} identifying the group of nuclear spins coupled to one another to form a total nuclear spin angular momentum.
\subsubsection*{\underline{Description}}
$I$ is the quantum number associated with the total nuclear spin angular momentum: $\boldsymbol{I} = \boldsymbol{I_1} + \boldsymbol{I_2} + \cdots$ where nuclei $1, 2, \cdots$ have individual nuclear spin angular momenta $\boldsymbol{I_1}, \boldsymbol{I_2}, \cdots$.
\subsubsection*{\underline{Restrictions}}
\begin{itemize}
\item non-negative integer or half-integer.
\end{itemize}
}
}
\vspace{0.5cm}
\fbox{
\parbox{12cm}{
\subsection*{$F_j$}
\subsubsection*{\underline{Element}}
\texttt{<lpcs:Fj>}
\subsubsection*{\underline{Attributes}}
\begin{itemize}\item\texttt{nuclearSpinRef}: a label identifying the nuclear spin being coupled to $\boldsymbol{J}$ or $\boldsymbol{F_{j-1}}$ to form an intermediate angular momentum;\item $j$: an integer label identifying the order of the hyperfine coupling\end{itemize}.
\subsubsection*{\underline{Description}}
$F_j$ is the intermediate angular momentum quantum number associated with the coupling of the nuclear spin angular momentum of nucleus $j$ to the intermediate angular momentum: $\boldsymbol{F_1} = \boldsymbol{J} + \boldsymbol{I_1}$ or $\boldsymbol{F_j} = \boldsymbol{F_{j-1}} + \boldsymbol{I_j}$; $\boldsymbol{F_j}$ is often not a good quantum number.
\subsubsection*{\underline{Restrictions}}
\begin{itemize}
\item non-negative integer or half-integer.
\end{itemize}
}
}
\vspace{0.5cm}
\fbox{
\parbox{12cm}{
\subsection*{$F$}
\subsubsection*{\underline{Element}}
\texttt{<lpcs:F>}
\subsubsection*{\underline{Attributes}}
\texttt{nuclearSpinRef}: a label, matching \texttt{/Q.+/} identifying the nuclear spin being coupled to $\boldsymbol{J}$ or $\boldsymbol{F_j}$ to form the total angular momentum.
\subsubsection*{\underline{Description}}
$F$ is the quantum number associated with the total angular momentum including nuclear spin: $\boldsymbol{F} = \boldsymbol{J} + \boldsymbol{I_1}$ if only one such coupling is resolved, $\boldsymbol{F} = \boldsymbol{F_{j-1}} + \boldsymbol{I_j}$ if two or more such couplings are resolved.
\subsubsection*{\underline{Restrictions}}
\begin{itemize}
\item non-negative integer or half-integer.
\end{itemize}
}
}
\vspace{0.5cm}
\fbox{
\parbox{12cm}{
\subsection*{$r$}
\subsubsection*{\underline{Element}}
\texttt{<lpcs:r>}
\subsubsection*{\underline{Attributes}}
\texttt{name}: a string identifying this ranking index.
\subsubsection*{\underline{Description}}
$r$ is a named, positive integer label identifying the state if no other good quantum numbers or symmetries are known.
\subsubsection*{\underline{Restrictions}}
\begin{itemize}
\item positive integer.
\end{itemize}
}
}
\vspace{0.5cm}
\fbox{
\parbox{12cm}{
\subsection*{\texttt{parity}}
\subsubsection*{\underline{Element}}
\texttt{<lpcs:parity>}
\subsubsection*{\underline{Attributes}}
None.
\subsubsection*{\underline{Description}}
\texttt{parity} is the total parity: the parity of the total molecular wavefunction (excluding nuclear spin) with respect to inversion through the molecular centre of mass of all particles' coordinates in the laboratory coordinate system, the $\hat{E}^*$ operation.
\subsubsection*{\underline{Restrictions}}
\begin{itemize}
\item `$+$' or `$-$'.
\end{itemize}
}
}
\vspace{0.5cm}
\fbox{
\parbox{12cm}{
\subsection*{\texttt{kronigParity}}
\subsubsection*{\underline{Element}}
\texttt{<lpcs:kronigParity>}
\subsubsection*{\underline{Attributes}}
None.
\subsubsection*{\underline{Description}}
\texttt{kronigParity} is the `rotationless' parity: the parity of the total molecular wavefunction excluding nuclear spin and rotation with respect to inversion through the molecular centre of mass of all particles' coordinates in the laboratory coordinate system. For integral $J$, the levels are called: \begin{align*}e\;\mathrm{if\;parity\;is}\;+(-1)^J,\\f\;\mathrm{if\;parity\;is}\;-(-1)^J.\end{align*}For half-odd integer $J$, the levels are called: \begin{align*}e\;\mathrm{if\;parity\;is}\;+(-1)^{J-\frac{1}{2}},\\f\;\mathrm{if\;parity\;is}\;-(-1)^{J-\frac{1}{2}}.\end{align*}.
\subsubsection*{\underline{Restrictions}}
\begin{itemize}
\item `e' or `f'.
\end{itemize}
}
}
\vspace{0.5cm}
\fbox{
\parbox{12cm}{
\subsection*{\texttt{asSym}}
\subsubsection*{\underline{Element}}
\texttt{<lpcs:asSym>}
\subsubsection*{\underline{Attributes}}
None.
\subsubsection*{\underline{Description}}
\texttt{asSym} is (for linear molecules with a centre of inversion) the symmetry of the rovibronic wavefunction: `a' or `s' such that the total wavefunction including nuclear spin is symmetric or antisymmetric with respect to permutation of the identical nuclei ($\hat{P}_{12}$), according to whether they are bosons or fermions respectively.
\subsubsection*{\underline{Restrictions}}
\begin{itemize}
\item `s' or `a'.
\end{itemize}
}
}
\vspace{0.5cm}




\section*{Closed-shell, asymmetric top molecules: asymcs}
\fbox{
\parbox{12cm}{
\subsection*{\texttt{ElecStateLabel}}
\subsubsection*{\underline{Element}}
\texttt{<asymcs:ElecStateLabel>}
\subsubsection*{\underline{Attributes}}
None.
\subsubsection*{\underline{Description}}
\texttt{ElecStateLabel} is a label identifying the electronic state.
\subsubsection*{\underline{Restrictions}}
\begin{itemize}
\item string.
\end{itemize}
}
}
\vspace{0.5cm}
\fbox{
\parbox{12cm}{
\subsection*{$v_i$}
\subsubsection*{\underline{Element}}
\texttt{<asymcs:vi>}
\subsubsection*{\underline{Attributes}}
\texttt{mode}: a positive integer, identifying the normal mode that this quantum number is associated with.
\subsubsection*{\underline{Description}}
$v_i$ is the vibrational quantum number associated with the $\nu_i$ normal mode.
\subsubsection*{\underline{Restrictions}}
\begin{itemize}
\item non-negative integer.
\end{itemize}
}
}
\vspace{0.5cm}
\fbox{
\parbox{12cm}{
\subsection*{\texttt{vibInv}}
\subsubsection*{\underline{Element}}
\texttt{<asymcs:vibInv>}
\subsubsection*{\underline{Attributes}}
None.
\subsubsection*{\underline{Description}}
\texttt{vibInv} is the parity of the vibrational wavefunction with respect to inversion through the molecular centre of mass in the molecular coordinate system..
\subsubsection*{\underline{Restrictions}}
\begin{itemize}
\item `s' or `a'.
\end{itemize}
}
}
\vspace{0.5cm}
\fbox{
\parbox{12cm}{
\subsection*{\texttt{vibSym}}
\subsubsection*{\underline{Element}}
\texttt{<asymcs:vibSym>}
\subsubsection*{\underline{Attributes}}
\texttt{group}: the symmetry group to which this species belongs.
\subsubsection*{\underline{Description}}
\texttt{vibSym} is the symmetry species of the vibrational wavefunction, in some appropriate symmetry group.
\subsubsection*{\underline{Restrictions}}
\begin{itemize}
\item string.
\end{itemize}
}
}
\vspace{0.5cm}
\fbox{
\parbox{12cm}{
\subsection*{$J$}
\subsubsection*{\underline{Element}}
\texttt{<asymcs:J>}
\subsubsection*{\underline{Attributes}}
None.
\subsubsection*{\underline{Description}}
$J$ is the quantum number associated with the total angular momentum excluding nuclear spin, $\boldsymbol{J}$.
\subsubsection*{\underline{Restrictions}}
\begin{itemize}
\item non-negative integer.
\end{itemize}
}
}
\vspace{0.5cm}
\fbox{
\parbox{12cm}{
\subsection*{$K_a$}
\subsubsection*{\underline{Element}}
\texttt{<asymcs:Ka>}
\subsubsection*{\underline{Attributes}}
None.
\subsubsection*{\underline{Description}}
$K_a$ is the rotational quantum label of an asymmetric top molecule, correlating to $K$ in the prolate symmetric top limit.
\subsubsection*{\underline{Restrictions}}
\begin{itemize}
\item non-negative integer.
\item $K_a \le J$.
\item $K_a + K_c = J \;\mathrm{or}\;J + 1$.
\end{itemize}
}
}
\vspace{0.5cm}
\fbox{
\parbox{12cm}{
\subsection*{$K_c$}
\subsubsection*{\underline{Element}}
\texttt{<asymcs:Kc>}
\subsubsection*{\underline{Attributes}}
None.
\subsubsection*{\underline{Description}}
$K_c$ is the rotational quantum label of an asymmetric top molecule, correlating to $K$ in the oblate symmetric top limit.
\subsubsection*{\underline{Restrictions}}
\begin{itemize}
\item non-negative integer.
\item $K_c \le J$.
\item $K_a + K_c = J \;\mathrm{or}\;J + 1$.
\end{itemize}
}
}
\vspace{0.5cm}
\fbox{
\parbox{12cm}{
\subsection*{\texttt{rotSym}}
\subsubsection*{\underline{Element}}
\texttt{<asymcs:rotSym>}
\subsubsection*{\underline{Attributes}}
\texttt{group}: the symmetry group to which this species belongs.
\subsubsection*{\underline{Description}}
\texttt{rotSym} is the symmetry species of the rotational wavefunction, in some appropriate symmetry group.
\subsubsection*{\underline{Restrictions}}
\begin{itemize}
\item string.
\end{itemize}
}
}
\vspace{0.5cm}
\fbox{
\parbox{12cm}{
\subsection*{$I$}
\subsubsection*{\underline{Element}}
\texttt{<asymcs:I>}
\subsubsection*{\underline{Attributes}}
\texttt{nuclearSpinRef}: a label identifying the group of nuclear spins coupled to one another to form a total nuclear spin angular momentum.
\subsubsection*{\underline{Description}}
$I$ is the quantum number associated with the total nuclear spin angular momentum: $\boldsymbol{I} = \boldsymbol{I_1} + \boldsymbol{I_2} + \cdots$ where nuclei $1, 2, \cdots$ have individual nuclear spin angular momenta $\boldsymbol{I_1}, \boldsymbol{I_2}, \cdots$.
\subsubsection*{\underline{Restrictions}}
\begin{itemize}
\item non-negative integer or half-integer.
\end{itemize}
}
}
\vspace{0.5cm}
\fbox{
\parbox{12cm}{
\subsection*{$F_j$}
\subsubsection*{\underline{Element}}
\texttt{<asymcs:Fj>}
\subsubsection*{\underline{Attributes}}
\begin{itemize}\item\texttt{nuclearSpinRef}: a label identifying the nuclear spin being coupled to $\boldsymbol{J}$ or $\boldsymbol{F_{j-1}}$ to form an intermediate angular momentum;\item $j$: an integer label identifying the order of the hyperfine coupling\end{itemize}.
\subsubsection*{\underline{Description}}
$F_j$ is the intermediate angular momentum quantum number associated with the coupling of the nuclear spin angular momentum of nucleus $j$ to the intermediate angular momentum: $\boldsymbol{F_1} = \boldsymbol{J} + \boldsymbol{I_1}$ or $\boldsymbol{F_j} = \boldsymbol{F_{j-1}} + \boldsymbol{I_j}$; $\boldsymbol{F_j}$ is often not a good quantum number.
\subsubsection*{\underline{Restrictions}}
\begin{itemize}
\item non-negative integer or half-integer.
\end{itemize}
}
}
\vspace{0.5cm}
\fbox{
\parbox{12cm}{
\subsection*{$F$}
\subsubsection*{\underline{Element}}
\texttt{<asymcs:F>}
\subsubsection*{\underline{Attributes}}
\texttt{nuclearSpinRef}: a label identifying the nuclear spin being coupled to $\boldsymbol{J}$ or $\boldsymbol{F_j}$ to form the total angular momentum.
\subsubsection*{\underline{Description}}
$F$ is the quantum number associated with the total angular momentum including nuclear spin: $\boldsymbol{F} = \boldsymbol{J} + \boldsymbol{I_1}$ if only one such coupling is resolved, $\boldsymbol{F} = \boldsymbol{F_{j-1}} + \boldsymbol{I_j}$ if two or more such couplings are resolved.
\subsubsection*{\underline{Restrictions}}
\begin{itemize}
\item non-negative integer or half-integer.
\end{itemize}
}
}
\vspace{0.5cm}
\fbox{
\parbox{12cm}{
\subsection*{$r$}
\subsubsection*{\underline{Element}}
\texttt{<asymcs:r>}
\subsubsection*{\underline{Attributes}}
\texttt{name}: a string identifying this ranking index.
\subsubsection*{\underline{Description}}
$r$ is a named, positive integer label identifying the state if no other good quantum numbers or symmetries are known.
\subsubsection*{\underline{Restrictions}}
\begin{itemize}
\item positive integer.
\end{itemize}
}
}
\vspace{0.5cm}
\fbox{
\parbox{12cm}{
\subsection*{\texttt{parity}}
\subsubsection*{\underline{Element}}
\texttt{<asymcs:parity>}
\subsubsection*{\underline{Attributes}}
None.
\subsubsection*{\underline{Description}}
\texttt{parity} is the total parity: the parity of the total molecular wavefunction (excluding nuclear spin) with respect to inversion through the molecular centre of mass of all particles' coordinates in the laboratory coordinate system, the $\hat{E}^*$ operation.
\subsubsection*{\underline{Restrictions}}
\begin{itemize}
\item `$+$' or `$-$'.
\end{itemize}
}
}
\vspace{0.5cm}




\section*{Open-shell, asymmetric top molecules: asymos}
\fbox{
\parbox{12cm}{
\subsection*{\texttt{ElecStateLabel}}
\subsubsection*{\underline{Element}}
\texttt{<asymos:ElecStateLabel>}
\subsubsection*{\underline{Attributes}}
None.
\subsubsection*{\underline{Description}}
\texttt{ElecStateLabel} is a label identifying the electronic state.
\subsubsection*{\underline{Restrictions}}
\begin{itemize}
\item string.
\end{itemize}
}
}
\vspace{0.5cm}
\fbox{
\parbox{12cm}{
\subsection*{\texttt{elecSym}}
\subsubsection*{\underline{Element}}
\texttt{<asymos:elecSym>}
\subsubsection*{\underline{Attributes}}
\texttt{group}: the symmetry group to which this species belongs.
\subsubsection*{\underline{Description}}
\texttt{elecSym} is the symmetry species of the electronic wavefunction described by some symmetry group..
\subsubsection*{\underline{Restrictions}}
\begin{itemize}
\item string.
\end{itemize}
}
}
\vspace{0.5cm}
\fbox{
\parbox{12cm}{
\subsection*{\texttt{elecInv}}
\subsubsection*{\underline{Element}}
\texttt{<asymos:elecInv>}
\subsubsection*{\underline{Attributes}}
None.
\subsubsection*{\underline{Description}}
\texttt{elecInv} is the parity of the electronic wavefunction with respect to inversion through the molecular centre of mass in the molecular coordinate system..
\subsubsection*{\underline{Restrictions}}
\begin{itemize}
\item `g' or `u'.
\end{itemize}
}
}
\vspace{0.5cm}
\fbox{
\parbox{12cm}{
\subsection*{$S$}
\subsubsection*{\underline{Element}}
\texttt{<asymos:S>}
\subsubsection*{\underline{Attributes}}
None.
\subsubsection*{\underline{Description}}
$S$ is the quantum number associated with the total electronic spin angular momentum, $\boldsymbol{S}$.
\subsubsection*{\underline{Restrictions}}
\begin{itemize}
\item non-negative integer or half-odd integer.
\end{itemize}
}
}
\vspace{0.5cm}
\fbox{
\parbox{12cm}{
\subsection*{$v_i$}
\subsubsection*{\underline{Element}}
\texttt{<asymos:vi>}
\subsubsection*{\underline{Attributes}}
\texttt{mode}: a positive integer, identifying the normal mode that this quantum number is associated with.
\subsubsection*{\underline{Description}}
$v_i$ is the vibrational quantum number associated with the $\nu_i$ normal mode.
\subsubsection*{\underline{Restrictions}}
\begin{itemize}
\item non-negative integer.
\end{itemize}
}
}
\vspace{0.5cm}
\fbox{
\parbox{12cm}{
\subsection*{\texttt{vibInv}}
\subsubsection*{\underline{Element}}
\texttt{<asymos:vibInv>}
\subsubsection*{\underline{Attributes}}
None.
\subsubsection*{\underline{Description}}
\texttt{vibInv} is the parity of the vibrational wavefunction with respect to inversion through the molecular centre of mass in the molecular coordinate system..
\subsubsection*{\underline{Restrictions}}
\begin{itemize}
\item `s' or `a'.
\end{itemize}
}
}
\vspace{0.5cm}
\fbox{
\parbox{12cm}{
\subsection*{\texttt{vibSym}}
\subsubsection*{\underline{Element}}
\texttt{<asymos:vibSym>}
\subsubsection*{\underline{Attributes}}
\texttt{group}: the symmetry group to which this species belongs.
\subsubsection*{\underline{Description}}
\texttt{vibSym} is the symmetry species of the vibrational wavefunction, in some appropriate symmetry group.
\subsubsection*{\underline{Restrictions}}
\begin{itemize}
\item string.
\end{itemize}
}
}
\vspace{0.5cm}
\fbox{
\parbox{12cm}{
\subsection*{$J$}
\subsubsection*{\underline{Element}}
\texttt{<asymos:J>}
\subsubsection*{\underline{Attributes}}
None.
\subsubsection*{\underline{Description}}
$J$ is the quantum number associated with the total angular momentum excluding nuclear spin, $\boldsymbol{J}$.
\subsubsection*{\underline{Restrictions}}
\begin{itemize}
\item non-negative integer or half-odd integer.
\end{itemize}
}
}
\vspace{0.5cm}
\fbox{
\parbox{12cm}{
\subsection*{$N$}
\subsubsection*{\underline{Element}}
\texttt{<asymos:N>}
\subsubsection*{\underline{Attributes}}
None.
\subsubsection*{\underline{Description}}
$N$ is the quantum number associated with the total angular momentum excluding electronic and nuclear spin, N: $\boldsymbol{J} = \boldsymbol{N} + \boldsymbol{S}$.
\subsubsection*{\underline{Restrictions}}
\begin{itemize}
\item non-negative integer.
\end{itemize}
}
}
\vspace{0.5cm}
\fbox{
\parbox{12cm}{
\subsection*{$K_a$}
\subsubsection*{\underline{Element}}
\texttt{<asymos:Ka>}
\subsubsection*{\underline{Attributes}}
None.
\subsubsection*{\underline{Description}}
$K_a$ is the rotational quantum label of an asymmetric top molecule, correlating to $K$ in the prolate symmetric top limit.
\subsubsection*{\underline{Restrictions}}
\begin{itemize}
\item non-negative integer.
\item $K_a \le N$.
\item $K_a + K_c = N \;\mathrm{or}\;N + 1$.
\end{itemize}
}
}
\vspace{0.5cm}
\fbox{
\parbox{12cm}{
\subsection*{$K_c$}
\subsubsection*{\underline{Element}}
\texttt{<asymos:Kc>}
\subsubsection*{\underline{Attributes}}
None.
\subsubsection*{\underline{Description}}
$K_c$ is the rotational quantum label of an asymmetric top molecule, correlating to $K$ in the oblate symmetric top limit.
\subsubsection*{\underline{Restrictions}}
\begin{itemize}
\item non-negative integer.
\item $K_c \le N$.
\item $K_a + K_c = N \;\mathrm{or}\;N + 1$.
\end{itemize}
}
}
\vspace{0.5cm}
\fbox{
\parbox{12cm}{
\subsection*{\texttt{rotSym}}
\subsubsection*{\underline{Element}}
\texttt{<asymos:rotSym>}
\subsubsection*{\underline{Attributes}}
\texttt{group}: the symmetry group to which this species belongs.
\subsubsection*{\underline{Description}}
\texttt{rotSym} is the symmetry species of the rotational wavefunction, in some appropriate symmetry group.
\subsubsection*{\underline{Restrictions}}
\begin{itemize}
\item string.
\end{itemize}
}
}
\vspace{0.5cm}
\fbox{
\parbox{12cm}{
\subsection*{$I$}
\subsubsection*{\underline{Element}}
\texttt{<asymos:I>}
\subsubsection*{\underline{Attributes}}
\texttt{nuclearSpinRef}: a label identifying the group of nuclear spins coupled to one another to form a total nuclear spin angular momentum.
\subsubsection*{\underline{Description}}
$I$ is the quantum number associated with the total nuclear spin angular momentum: $\boldsymbol{I} = \boldsymbol{I_1} + \boldsymbol{I_2} + \cdots$ where nuclei $1, 2, \cdots$ have individual nuclear spin angular momenta $\boldsymbol{I_1}, \boldsymbol{I_2}, \cdots$.
\subsubsection*{\underline{Restrictions}}
\begin{itemize}
\item non-negative integer or half-integer.
\end{itemize}
}
}
\vspace{0.5cm}
\fbox{
\parbox{12cm}{
\subsection*{$F_j$}
\subsubsection*{\underline{Element}}
\texttt{<asymos:Fj>}
\subsubsection*{\underline{Attributes}}
\begin{itemize}\item\texttt{nuclearSpinRef}: a label identifying the nuclear spin being coupled to $\boldsymbol{J}$ or $\boldsymbol{F_{j-1}}$ to form an intermediate angular momentum;\item $j$: an integer label identifying the order of the hyperfine coupling\end{itemize}.
\subsubsection*{\underline{Description}}
$F_j$ is the intermediate angular momentum quantum number associated with the coupling of the nuclear spin angular momentum of nucleus $j$ to the intermediate angular momentum: $\boldsymbol{F_1} = \boldsymbol{J} + \boldsymbol{I_1}$ or $\boldsymbol{F_j} = \boldsymbol{F_{j-1}} + \boldsymbol{I_j}$; $\boldsymbol{F_j}$ is often not a good quantum number.
\subsubsection*{\underline{Restrictions}}
\begin{itemize}
\item non-negative integer or half-integer.
\end{itemize}
}
}
\vspace{0.5cm}
\fbox{
\parbox{12cm}{
\subsection*{$F$}
\subsubsection*{\underline{Element}}
\texttt{<asymos:F>}
\subsubsection*{\underline{Attributes}}
\texttt{nuclearSpinRef}: a label identifying the nuclear spin being coupled to $\boldsymbol{J}$ or $\boldsymbol{F_j}$ to form the total angular momentum.
\subsubsection*{\underline{Description}}
$F$ is the quantum number associated with the total angular momentum including nuclear spin: $\boldsymbol{F} = \boldsymbol{J} + \boldsymbol{I_1}$ if only one such coupling is resolved, $\boldsymbol{F} = \boldsymbol{F_{j-1}} + \boldsymbol{I_j}$ if two or more such couplings are resolved.
\subsubsection*{\underline{Restrictions}}
\begin{itemize}
\item non-negative integer or half-integer.
\end{itemize}
}
}
\vspace{0.5cm}
\fbox{
\parbox{12cm}{
\subsection*{$r$}
\subsubsection*{\underline{Element}}
\texttt{<asymos:r>}
\subsubsection*{\underline{Attributes}}
\texttt{name}: a string identifying this ranking index.
\subsubsection*{\underline{Description}}
$r$ is a named, positive integer label identifying the state if no other good quantum numbers or symmetries are known.
\subsubsection*{\underline{Restrictions}}
\begin{itemize}
\item positive integer.
\end{itemize}
}
}
\vspace{0.5cm}
\fbox{
\parbox{12cm}{
\subsection*{\texttt{parity}}
\subsubsection*{\underline{Element}}
\texttt{<asymos:parity>}
\subsubsection*{\underline{Attributes}}
None.
\subsubsection*{\underline{Description}}
\texttt{parity} is the total parity: the parity of the total molecular wavefunction (excluding nuclear spin) with respect to inversion through the molecular centre of mass of all particles' coordinates in the laboratory coordinate system, the $\hat{E}^*$ operation.
\subsubsection*{\underline{Restrictions}}
\begin{itemize}
\item `$+$' or `$-$'.
\end{itemize}
}
}
\vspace{0.5cm}




\section*{Closed-shell, spherical-top molecules: sphcs}
\fbox{
\parbox{12cm}{
\subsection*{\texttt{ElecStateLabel}}
\subsubsection*{\underline{Element}}
\texttt{<sphcs:ElecStateLabel>}
\subsubsection*{\underline{Attributes}}
None.
\subsubsection*{\underline{Description}}
\texttt{ElecStateLabel} is a label identifying the electronic state.
\subsubsection*{\underline{Restrictions}}
\begin{itemize}
\item string.
\end{itemize}
}
}
\vspace{0.5cm}
\fbox{
\parbox{12cm}{
\subsection*{$v_i$}
\subsubsection*{\underline{Element}}
\texttt{<sphcs:vi>}
\subsubsection*{\underline{Attributes}}
\texttt{mode}: a positive integer, identifying the normal mode that this quantum number is associated with.
\subsubsection*{\underline{Description}}
$v_i$ is the vibrational quantum number associated with the $\nu_i$ normal mode.
\subsubsection*{\underline{Restrictions}}
\begin{itemize}
\item non-negative integer.
\end{itemize}
}
}
\vspace{0.5cm}
\fbox{
\parbox{12cm}{
\subsection*{$l_i$}
\subsubsection*{\underline{Element}}
\texttt{<sphcs:li>}
\subsubsection*{\underline{Attributes}}
\texttt{mode}: a positive integer, identifying the degenerate normal mode that this vibrational angular momentum quantum number is associated with.
\subsubsection*{\underline{Description}}
$l_i$ is the vibrational angular momentum quantum number associated with the degenerate $\nu_i$ normal mode; positive and negative values distinguish $l$-type doubling components.
\subsubsection*{\underline{Restrictions}}
\begin{itemize}
\item non-negative integer.
\item $|l_i| = v_i, v_i-2, \cdots, 1 \;\mathrm{or}\;0$.
\end{itemize}
}
}
\vspace{0.5cm}
\fbox{
\parbox{12cm}{
\subsection*{\texttt{vibSym}}
\subsubsection*{\underline{Element}}
\texttt{<sphcs:vibSym>}
\subsubsection*{\underline{Attributes}}
\texttt{group}: the symmetry group to which this species belongs.
\subsubsection*{\underline{Description}}
\texttt{vibSym} is the symmetry species of the vibrational wavefunction, in some appropriate symmetry group.
\subsubsection*{\underline{Restrictions}}
\begin{itemize}
\item string.
\end{itemize}
}
}
\vspace{0.5cm}
\fbox{
\parbox{12cm}{
\subsection*{$J$}
\subsubsection*{\underline{Element}}
\texttt{<sphcs:J>}
\subsubsection*{\underline{Attributes}}
\texttt{group}: the symmetry group to which this species belongs.
\subsubsection*{\underline{Description}}
$J$ is the quantum number associated with the total angular momentum excluding nuclear spin, $\boldsymbol{J}$.
\subsubsection*{\underline{Restrictions}}
\begin{itemize}
\item non-negative integer.
\end{itemize}
}
}
\vspace{0.5cm}
\fbox{
\parbox{12cm}{
\subsection*{\texttt{rotSym}}
\subsubsection*{\underline{Element}}
\texttt{<sphcs:rotSym>}
\subsubsection*{\underline{Attributes}}
None.
\subsubsection*{\underline{Description}}
\texttt{rotSym} is the symmetry species of the rotational wavefunction, in some appropriate symmetry group.
\subsubsection*{\underline{Restrictions}}
\begin{itemize}
\item string.
\end{itemize}
}
}
\vspace{0.5cm}
\fbox{
\parbox{12cm}{
\subsection*{$I$}
\subsubsection*{\underline{Element}}
\texttt{<sphcs:I>}
\subsubsection*{\underline{Attributes}}
\texttt{nuclearSpinRef}: a label identifying the group of nuclear spins coupled to one another to form a total nuclear spin angular momentum.
\subsubsection*{\underline{Description}}
$I$ is the quantum number associated with the total nuclear spin angular momentum: $\boldsymbol{I} = \boldsymbol{I_1} + \boldsymbol{I_2} + \cdots$ where nuclei $1, 2, \cdots$ have individual nuclear spin angular momenta $\boldsymbol{I_1}, \boldsymbol{I_2}, \cdots$.
\subsubsection*{\underline{Restrictions}}
\begin{itemize}
\item non-negative integer or half-integer.
\end{itemize}
}
}
\vspace{0.5cm}
\fbox{
\parbox{12cm}{
\subsection*{$F_j$}
\subsubsection*{\underline{Element}}
\texttt{<sphcs:Fj>}
\subsubsection*{\underline{Attributes}}
\begin{itemize}\item\texttt{nuclearSpinRef}: a label identifying the nuclear spin being coupled to $\boldsymbol{J}$ or $\boldsymbol{F_{j-1}}$ to form an intermediate angular momentum;\item $j$: an integer label identifying the order of the hyperfine coupling\end{itemize}.
\subsubsection*{\underline{Description}}
$F_j$ is the intermediate angular momentum quantum number associated with the coupling of the nuclear spin angular momentum of nucleus $j$ to the intermediate angular momentum: $\boldsymbol{F_1} = \boldsymbol{J} + \boldsymbol{I_1}$ or $\boldsymbol{F_j} = \boldsymbol{F_{j-1}} + \boldsymbol{I_j}$; $\boldsymbol{F_j}$ is often not a good quantum number.
\subsubsection*{\underline{Restrictions}}
\begin{itemize}
\item non-negative integer or half-integer.
\end{itemize}
}
}
\vspace{0.5cm}
\fbox{
\parbox{12cm}{
\subsection*{$F$}
\subsubsection*{\underline{Element}}
\texttt{<sphcs:F>}
\subsubsection*{\underline{Attributes}}
\texttt{nuclearSpinRef}: a label identifying the nuclear spin being coupled to $\boldsymbol{J}$ or $\boldsymbol{F_j}$ to form the total angular momentum.
\subsubsection*{\underline{Description}}
$F$ is the quantum number associated with the total angular momentum including nuclear spin: $\boldsymbol{F} = \boldsymbol{J} + \boldsymbol{I_1}$ if only one such coupling is resolved, $\boldsymbol{F} = \boldsymbol{F_{j-1}} + \boldsymbol{I_j}$ if two or more such couplings are resolved.
\subsubsection*{\underline{Restrictions}}
\begin{itemize}
\item non-negative integer or half-integer.
\end{itemize}
}
}
\vspace{0.5cm}
\fbox{
\parbox{12cm}{
\subsection*{$r$}
\subsubsection*{\underline{Element}}
\texttt{<sphcs:r>}
\subsubsection*{\underline{Attributes}}
\texttt{name}: a string identifying this ranking index.
\subsubsection*{\underline{Description}}
$r$ is a named, positive integer label identifying the state if no other good quantum numbers or symmetries are known.
\subsubsection*{\underline{Restrictions}}
\begin{itemize}
\item positive integer.
\end{itemize}
}
}
\vspace{0.5cm}
\fbox{
\parbox{12cm}{
\subsection*{\texttt{parity}}
\subsubsection*{\underline{Element}}
\texttt{<sphcs:parity>}
\subsubsection*{\underline{Attributes}}
None.
\subsubsection*{\underline{Description}}
\texttt{parity} is the total parity: the parity of the total molecular wavefunction (excluding nuclear spin) with respect to inversion through the molecular centre of mass of all particles' coordinates in the laboratory coordinate system, the $\hat{E}^*$ operation.
\subsubsection*{\underline{Restrictions}}
\begin{itemize}
\item `$+$' or `$-$'.
\end{itemize}
}
}
\vspace{0.5cm}




\section*{Open-shell, spherical-top molecules: sphos}
\fbox{
\parbox{12cm}{
\subsection*{\texttt{ElecStateLabel}}
\subsubsection*{\underline{Element}}
\texttt{<sphos:ElecStateLabel>}
\subsubsection*{\underline{Attributes}}
None.
\subsubsection*{\underline{Description}}
\texttt{ElecStateLabel} is a label identifying the electronic state.
\subsubsection*{\underline{Restrictions}}
\begin{itemize}
\item string.
\end{itemize}
}
}
\vspace{0.5cm}
\fbox{
\parbox{12cm}{
\subsection*{\texttt{elecSym}}
\subsubsection*{\underline{Element}}
\texttt{<sphos:elecSym>}
\subsubsection*{\underline{Attributes}}
\texttt{group}: the symmetry group to which this species belongs.
\subsubsection*{\underline{Description}}
\texttt{elecSym} is the symmetry species of the electronic wavefunction described by some symmetry group..
\subsubsection*{\underline{Restrictions}}
\begin{itemize}
\item string.
\end{itemize}
}
}
\vspace{0.5cm}
\fbox{
\parbox{12cm}{
\subsection*{\texttt{elecInv}}
\subsubsection*{\underline{Element}}
\texttt{<sphos:elecInv>}
\subsubsection*{\underline{Attributes}}
None.
\subsubsection*{\underline{Description}}
\texttt{elecInv} is the parity of the electronic wavefunction with respect to inversion through the molecular centre of mass in the molecular coordinate system..
\subsubsection*{\underline{Restrictions}}
\begin{itemize}
\item `g' or `u'.
\end{itemize}
}
}
\vspace{0.5cm}
\fbox{
\parbox{12cm}{
\subsection*{$S$}
\subsubsection*{\underline{Element}}
\texttt{<sphos:S>}
\subsubsection*{\underline{Attributes}}
None.
\subsubsection*{\underline{Description}}
$S$ is the quantum number associated with the total electronic spin angular momentum, $\boldsymbol{S}$.
\subsubsection*{\underline{Restrictions}}
\begin{itemize}
\item non-negative integer or half-odd integer.
\end{itemize}
}
}
\vspace{0.5cm}
\fbox{
\parbox{12cm}{
\subsection*{$v_i$}
\subsubsection*{\underline{Element}}
\texttt{<sphos:vi>}
\subsubsection*{\underline{Attributes}}
\texttt{mode}: a positive integer, identifying the normal mode that this quantum number is associated with.
\subsubsection*{\underline{Description}}
$v_i$ is the vibrational quantum number associated with the $\nu_i$ normal mode.
\subsubsection*{\underline{Restrictions}}
\begin{itemize}
\item non-negative integer.
\end{itemize}
}
}
\vspace{0.5cm}
\fbox{
\parbox{12cm}{
\subsection*{$l_i$}
\subsubsection*{\underline{Element}}
\texttt{<sphos:li>}
\subsubsection*{\underline{Attributes}}
\texttt{mode}: a positive integer, identifying the degenerate normal mode that this vibrational angular momentum quantum number is associated with.
\subsubsection*{\underline{Description}}
$l_i$ is the vibrational angular momentum quantum number associated with the degenerate $\nu_i$ normal mode; positive and negative values distinguish $l$-type doubling components.
\subsubsection*{\underline{Restrictions}}
\begin{itemize}
\item non-negative integer.
\item $|l_i| = v_i, v_i-2, \cdots, 1 \;\mathrm{or}\;0$.
\end{itemize}
}
}
\vspace{0.5cm}
\fbox{
\parbox{12cm}{
\subsection*{\texttt{vibSym}}
\subsubsection*{\underline{Element}}
\texttt{<sphos:vibSym>}
\subsubsection*{\underline{Attributes}}
\texttt{group}: the symmetry group to which this species belongs.
\subsubsection*{\underline{Description}}
\texttt{vibSym} is the symmetry species of the vibrational wavefunction, in some appropriate symmetry group.
\subsubsection*{\underline{Restrictions}}
\begin{itemize}
\item string.
\end{itemize}
}
}
\vspace{0.5cm}
\fbox{
\parbox{12cm}{
\subsection*{$J$}
\subsubsection*{\underline{Element}}
\texttt{<sphos:J>}
\subsubsection*{\underline{Attributes}}
None.
\subsubsection*{\underline{Description}}
$J$ is the quantum number associated with the total angular momentum excluding nuclear spin, $\boldsymbol{J}$.
\subsubsection*{\underline{Restrictions}}
\begin{itemize}
\item non-negative integer or half-odd integer.
\end{itemize}
}
}
\vspace{0.5cm}
\fbox{
\parbox{12cm}{
\subsection*{$N$}
\subsubsection*{\underline{Element}}
\texttt{<sphos:N>}
\subsubsection*{\underline{Attributes}}
None.
\subsubsection*{\underline{Description}}
$N$ is the quantum number associated with the total angular momentum excluding electronic and nuclear spin, N: $\boldsymbol{J} = \boldsymbol{N} + \boldsymbol{S}$.
\subsubsection*{\underline{Restrictions}}
\begin{itemize}
\item non-negative integer.
\end{itemize}
}
}
\vspace{0.5cm}
\fbox{
\parbox{12cm}{
\subsection*{\texttt{rotSym}}
\subsubsection*{\underline{Element}}
\texttt{<sphos:rotSym>}
\subsubsection*{\underline{Attributes}}
\texttt{group}: the symmetry group to which this species belongs.
\subsubsection*{\underline{Description}}
\texttt{rotSym} is the symmetry species of the rotational wavefunction, in some appropriate symmetry group.
\subsubsection*{\underline{Restrictions}}
\begin{itemize}
\item string.
\end{itemize}
}
}
\vspace{0.5cm}
\fbox{
\parbox{12cm}{
\subsection*{$I$}
\subsubsection*{\underline{Element}}
\texttt{<sphos:I>}
\subsubsection*{\underline{Attributes}}
\texttt{nuclearSpinRef}: a label identifying the group of nuclear spins coupled to one another to form a total nuclear spin angular momentum.
\subsubsection*{\underline{Description}}
$I$ is the quantum number associated with the total nuclear spin angular momentum: $\boldsymbol{I} = \boldsymbol{I_1} + \boldsymbol{I_2} + \cdots$ where nuclei $1, 2, \cdots$ have individual nuclear spin angular momenta $\boldsymbol{I_1}, \boldsymbol{I_2}, \cdots$.
\subsubsection*{\underline{Restrictions}}
\begin{itemize}
\item non-negative integer or half-integer.
\end{itemize}
}
}
\vspace{0.5cm}
\fbox{
\parbox{12cm}{
\subsection*{$F_j$}
\subsubsection*{\underline{Element}}
\texttt{<sphos:Fj>}
\subsubsection*{\underline{Attributes}}
\begin{itemize}\item\texttt{nuclearSpinRef}: a label identifying the nuclear spin being coupled to $\boldsymbol{J}$ or $\boldsymbol{F_{j-1}}$ to form an intermediate angular momentum;\item $j$: an integer label identifying the order of the hyperfine coupling\end{itemize}.
\subsubsection*{\underline{Description}}
$F_j$ is the intermediate angular momentum quantum number associated with the coupling of the nuclear spin angular momentum of nucleus $j$ to the intermediate angular momentum: $\boldsymbol{F_1} = \boldsymbol{J} + \boldsymbol{I_1}$ or $\boldsymbol{F_j} = \boldsymbol{F_{j-1}} + \boldsymbol{I_j}$; $\boldsymbol{F_j}$ is often not a good quantum number.
\subsubsection*{\underline{Restrictions}}
\begin{itemize}
\item non-negative integer or half-integer.
\end{itemize}
}
}
\vspace{0.5cm}
\fbox{
\parbox{12cm}{
\subsection*{$F$}
\subsubsection*{\underline{Element}}
\texttt{<sphos:F>}
\subsubsection*{\underline{Attributes}}
\texttt{nuclearSpinRef}: a label identifying the nuclear spin being coupled to $\boldsymbol{J}$ or $\boldsymbol{F_j}$ to form the total angular momentum.
\subsubsection*{\underline{Description}}
$F$ is the quantum number associated with the total angular momentum including nuclear spin: $\boldsymbol{F} = \boldsymbol{J} + \boldsymbol{I_1}$ if only one such coupling is resolved, $\boldsymbol{F} = \boldsymbol{F_{j-1}} + \boldsymbol{I_j}$ if two or more such couplings are resolved.
\subsubsection*{\underline{Restrictions}}
\begin{itemize}
\item non-negative integer or half-integer.
\end{itemize}
}
}
\vspace{0.5cm}
\fbox{
\parbox{12cm}{
\subsection*{$r$}
\subsubsection*{\underline{Element}}
\texttt{<sphos:r>}
\subsubsection*{\underline{Attributes}}
\texttt{name}: a string identifying this ranking index.
\subsubsection*{\underline{Description}}
$r$ is a named, positive integer label identifying the state if no other good quantum numbers or symmetries are known.
\subsubsection*{\underline{Restrictions}}
\begin{itemize}
\item positive integer.
\end{itemize}
}
}
\vspace{0.5cm}
\fbox{
\parbox{12cm}{
\subsection*{\texttt{parity}}
\subsubsection*{\underline{Element}}
\texttt{<sphos:parity>}
\subsubsection*{\underline{Attributes}}
None.
\subsubsection*{\underline{Description}}
\texttt{parity} is the total parity: the parity of the total molecular wavefunction (excluding nuclear spin) with respect to inversion through the molecular centre of mass of all particles' coordinates in the laboratory coordinate system, the $\hat{E}^*$ operation.
\subsubsection*{\underline{Restrictions}}
\begin{itemize}
\item `$+$' or `$-$'.
\end{itemize}
}
}
\vspace{0.5cm}




\section*{Open-shell,linear triatomic molecules: ltos}
\fbox{
\parbox{12cm}{
\subsection*{\texttt{ElecStateLabel}}
\subsubsection*{\underline{Element}}
\texttt{<ltos:ElecStateLabel>}
\subsubsection*{\underline{Attributes}}
None.
\subsubsection*{\underline{Description}}
\texttt{ElecStateLabel} is a label identifying the electronic state: $X$, $A$, $a$, $B$, etc..
\subsubsection*{\underline{Restrictions}}
\begin{itemize}
\item string.
\end{itemize}
}
}
\vspace{0.5cm}
\fbox{
\parbox{12cm}{
\subsection*{\texttt{elecInv}}
\subsubsection*{\underline{Element}}
\texttt{<ltos:elecInv>}
\subsubsection*{\underline{Attributes}}
None.
\subsubsection*{\underline{Description}}
\texttt{elecInv} is the parity of the electronic wavefunction with respect to inversion through the molecular centre of mass in the molecular coordinate system.
\subsubsection*{\underline{Restrictions}}
\begin{itemize}
\item `g' or `u'.
\end{itemize}
}
}
\vspace{0.5cm}
\fbox{
\parbox{12cm}{
\subsection*{\texttt{elecRefl}}
\subsubsection*{\underline{Element}}
\texttt{<ltos:elecRefl>}
\subsubsection*{\underline{Attributes}}
None.
\subsubsection*{\underline{Description}}
\texttt{elecRefl} is the parity of the electronic wavefunction with respect to reflection in a plane containing the molecular symmetry axis in the molecular coordinate system (equivalent to inversion through the molecular centre of mass in the laboratory coordinate system).
\subsubsection*{\underline{Restrictions}}
\begin{itemize}
\item `$+$' or `$-$'.
\end{itemize}
}
}
\vspace{0.5cm}
\fbox{
\parbox{12cm}{
\subsection*{$S$}
\subsubsection*{\underline{Element}}
\texttt{<ltos:S>}
\subsubsection*{\underline{Attributes}}
None.
\subsubsection*{\underline{Description}}
$S$ is the quantum number associated with the total electronic spin angular momentum, $\boldsymbol{S}$.
\subsubsection*{\underline{Restrictions}}
\begin{itemize}
\item non-negative integer or half-odd integer.
\end{itemize}
}
}
\vspace{0.5cm}
\fbox{
\parbox{12cm}{
\subsection*{$N$}
\subsubsection*{\underline{Element}}
\texttt{<ltos:N>}
\subsubsection*{\underline{Attributes}}
None.
\subsubsection*{\underline{Description}}
$N$ is the quantum number associated with the total angular momentum excluding electronic and nuclear spin, N: $\boldsymbol{J} = \boldsymbol{N} + \boldsymbol{S}$.
\subsubsection*{\underline{Restrictions}}
\begin{itemize}
\item non-negative integer.
\end{itemize}
}
}
\vspace{0.5cm}
\fbox{
\parbox{12cm}{
\subsection*{$v_1$}
\subsubsection*{\underline{Element}}
\texttt{<ltos:v1>}
\subsubsection*{\underline{Attributes}}
None.
\subsubsection*{\underline{Description}}
$v_1$ is the vibrational quantum number associated with the $\nu_1$ normal mode.
\subsubsection*{\underline{Restrictions}}
\begin{itemize}
\item non-negative integer.
\end{itemize}
}
}
\vspace{0.5cm}
\fbox{
\parbox{12cm}{
\subsection*{$v_2$}
\subsubsection*{\underline{Element}}
\texttt{<ltos:v2>}
\subsubsection*{\underline{Attributes}}
None.
\subsubsection*{\underline{Description}}
$v_2$ is the vibrational quantum number associated with the doubly-degenerate $\nu_2$ normal mode.
\subsubsection*{\underline{Restrictions}}
\begin{itemize}
\item non-negative integer.
\end{itemize}
}
}
\vspace{0.5cm}
\fbox{
\parbox{12cm}{
\subsection*{$l_2$}
\subsubsection*{\underline{Element}}
\texttt{<ltos:l2>}
\subsubsection*{\underline{Attributes}}
None.
\subsubsection*{\underline{Description}}
$l_2$ is the vibrational angular momentum quantum number associated with the degenerate bending vibration, $\nu_2$; positive and negative values distinguish $l$-type doubling components.
\subsubsection*{\underline{Restrictions}}
\begin{itemize}
\item integer.
\item $|l_2| = v_2, v_2 - 2, \cdots, 1 \;\mathrm{or}\;0$.
\end{itemize}
}
}
\vspace{0.5cm}
\fbox{
\parbox{12cm}{
\subsection*{$v_3$}
\subsubsection*{\underline{Element}}
\texttt{<ltos:v3>}
\subsubsection*{\underline{Attributes}}
None.
\subsubsection*{\underline{Description}}
$v_3$ is the vibrational quantum number associated with the $\nu_3$ normal mode.
\subsubsection*{\underline{Restrictions}}
\begin{itemize}
\item non-negative integer.
\end{itemize}
}
}
\vspace{0.5cm}
\fbox{
\parbox{12cm}{
\subsection*{$J$}
\subsubsection*{\underline{Element}}
\texttt{<ltos:J>}
\subsubsection*{\underline{Attributes}}
None.
\subsubsection*{\underline{Description}}
$J$ is the quantum number associated with the total angular momentum excluding nuclear spin, $\boldsymbol{J}$.
\subsubsection*{\underline{Restrictions}}
\begin{itemize}
\item non-negative integer.
\end{itemize}
}
}
\vspace{0.5cm}
\fbox{
\parbox{12cm}{
\subsection*{$F_1$}
\subsubsection*{\underline{Element}}
\texttt{<ltos:F1>}
\subsubsection*{\underline{Attributes}}
\texttt{nuclearSpinRef}: a label identifying the nuclear spin coupled to $\boldsymbol{J}$ to form the intermediate angular momentum.
\subsubsection*{\underline{Description}}
$F_1$ is the intermediate angular momentum quantum number associated with the coupling of the rotational angular momentum and nuclear spin of nucleus 1 where two or more such couplings are resolved: $\boldsymbol{F_1} = \boldsymbol{J} + \boldsymbol{I_1}$; $F_1$ is often not a good quantum number.
\subsubsection*{\underline{Restrictions}}
\begin{itemize}
\item non-negative integer or half-integer.
\item $|J - I_1| \le F_1 \le J + I_1$.
\end{itemize}
}
}
\vspace{0.5cm}
\fbox{
\parbox{12cm}{
\subsection*{$F_2$}
\subsubsection*{\underline{Element}}
\texttt{<ltos:F2>}
\subsubsection*{\underline{Attributes}}
\texttt{nuclearSpinRef}: a label identifying the nuclear spin coupled to $\boldsymbol{F_1}$ to form an intermediate angular momentum.
\subsubsection*{\underline{Description}}
$F_2$ is the intermediate angular momentum quantum number associated with the coupling of the rotational angular momentum and nuclear spin of nucleus 2 where three such couplings are resolved: $\boldsymbol{F_2} = \boldsymbol{F_1} + \boldsymbol{I_2}$; $F_2$ is often not a good quantum number.
\subsubsection*{\underline{Restrictions}}
\begin{itemize}
\item non-negative integer or half-integer.
\item $|F_1 - I_2| \le F_2 \le F_1 + I_2$.
\end{itemize}
}
}
\vspace{0.5cm}
\fbox{
\parbox{12cm}{
\subsection*{$F$}
\subsubsection*{\underline{Element}}
\texttt{<ltos:F>}
\subsubsection*{\underline{Attributes}}
\texttt{nuclearSpinRef}: a label identifying the nuclear spin coupled to $\boldsymbol{J}$, $\boldsymbol{F_1}$, or $\boldsymbol{F}$ to form the total angular momentum.
\subsubsection*{\underline{Description}}
$F$ is the quantum number associated with the total angular momentum including nuclear spin: $\boldsymbol{F} = \boldsymbol{J} + \boldsymbol{I_1}$ if only one hyperfine coupling is resolved, $\boldsymbol{F} = \boldsymbol{F_1} + \boldsymbol{I_2}$ if two couplings are resolved, or $\boldsymbol{F} = \boldsymbol{F_2} + \boldsymbol{I_3}$ if all three couplings are resolved.
\subsubsection*{\underline{Restrictions}}
\begin{itemize}
\item non-negative integer or half-integer.
\item $|F_2 - I_3| \le F \le F_2 + I_3$.
\end{itemize}
}
}
\vspace{0.5cm}
\fbox{
\parbox{12cm}{
\subsection*{$r$}
\subsubsection*{\underline{Element}}
\texttt{<ltos:r>}
\subsubsection*{\underline{Attributes}}
\texttt{name}: a string identifying this ranking index.
\subsubsection*{\underline{Description}}
$r$ is a named, positive integer label identifying the state if no other good quantum numbers or symmetries are known.
\subsubsection*{\underline{Restrictions}}
\begin{itemize}
\item positive integer.
\end{itemize}
}
}
\vspace{0.5cm}
\fbox{
\parbox{12cm}{
\subsection*{\texttt{parity}}
\subsubsection*{\underline{Element}}
\texttt{<ltos:parity>}
\subsubsection*{\underline{Attributes}}
None.
\subsubsection*{\underline{Description}}
\texttt{parity} is the total parity: the parity of the total molecular wavefunction (excluding nuclear spin) with respect to inversion through the molecular centre of mass of all particles' coordinates in the laboratory coordinate system, the $\hat{E}^*$ operation.
\subsubsection*{\underline{Restrictions}}
\begin{itemize}
\item `$+$' or `$-$'.
\end{itemize}
}
}
\vspace{0.5cm}
\fbox{
\parbox{12cm}{
\subsection*{\texttt{kronigParity}}
\subsubsection*{\underline{Element}}
\texttt{<ltos:kronigParity>}
\subsubsection*{\underline{Attributes}}
None.
\subsubsection*{\underline{Description}}
\texttt{kronigParity} is the `rotationless' parity: the parity of the total molecular wavefunction excluding nuclear spin and rotation with respect to inversion through the molecular centre of mass of all particles' coordinates in the laboratory coordinate system. For integral $J$, the levels are called: \begin{align*}e\;\mathrm{if\;parity\;is}\;+(-1)^J,\\f\;\mathrm{if\;parity\;is}\;-(-1)^J.\end{align*}For half-odd integer $J$, the levels are called: \begin{align*}e\;\mathrm{if\;parity\;is}\;+(-1)^{J-\frac{1}{2}},\\f\;\mathrm{if\;parity\;is}\;-(-1)^{J-\frac{1}{2}}.\end{align*}.
\subsubsection*{\underline{Restrictions}}
\begin{itemize}
\item `e' or `f'.
\end{itemize}
}
}
\vspace{0.5cm}
\fbox{
\parbox{12cm}{
\subsection*{\texttt{asSym}}
\subsubsection*{\underline{Element}}
\texttt{<ltos:asSym>}
\subsubsection*{\underline{Attributes}}
None.
\subsubsection*{\underline{Description}}
\texttt{asSym} is (for linear molecules with a centre of inversion) the symmetry of the rovibronic wavefunction: `a' or `s' such that the total wavefunction including nuclear spin is symmetric or antisymmetric with respect to permutation of the identical nuclei ($\hat{P}_{12}$), according to whether they are bosons or fermions respectively.
\subsubsection*{\underline{Restrictions}}
\begin{itemize}
\item `s' or `a'.
\end{itemize}
}
}
\vspace{0.5cm}




\section*{Open-shell, linear polyatomic molecules: lpos}
\fbox{
\parbox{12cm}{
\subsection*{\texttt{ElecStateLabel}}
\subsubsection*{\underline{Element}}
\texttt{<lpos:ElecStateLabel>}
\subsubsection*{\underline{Attributes}}
None.
\subsubsection*{\underline{Description}}
\texttt{ElecStateLabel} is a label identifying the electronic state.
\subsubsection*{\underline{Restrictions}}
\begin{itemize}
\item string.
\end{itemize}
}
}
\vspace{0.5cm}
\fbox{
\parbox{12cm}{
\subsection*{\texttt{elecInv}}
\subsubsection*{\underline{Element}}
\texttt{<lpos:elecInv>}
\subsubsection*{\underline{Attributes}}
None.
\subsubsection*{\underline{Description}}
\texttt{elecInv} is the parity of the electronic wavefunction with respect to inversion through the molecular centre of mass in the molecular coordinate system.
\subsubsection*{\underline{Restrictions}}
\begin{itemize}
\item `g' or `u'.
\end{itemize}
}
}
\vspace{0.5cm}
\fbox{
\parbox{12cm}{
\subsection*{\texttt{elecRefl}}
\subsubsection*{\underline{Element}}
\texttt{<lpos:elecRefl>}
\subsubsection*{\underline{Attributes}}
None.
\subsubsection*{\underline{Description}}
\texttt{elecRefl} is the parity of the electronic wavefunction with respect to reflection in a plane containing the molecular symmetry axis in the molecular coordinate system (equivalent to inversion through the molecular centre of mass in the laboratory coordinate system).
\subsubsection*{\underline{Restrictions}}
\begin{itemize}
\item `$+$' or `$-$'.
\end{itemize}
}
}
\vspace{0.5cm}
\fbox{
\parbox{12cm}{
\subsection*{$S$}
\subsubsection*{\underline{Element}}
\texttt{<lpos:S>}
\subsubsection*{\underline{Attributes}}
None.
\subsubsection*{\underline{Description}}
$S$ is the quantum number associated with the total electronic spin angular momentum, $\boldsymbol{S}$.
\subsubsection*{\underline{Restrictions}}
\begin{itemize}
\item non-negative integer or half-odd integer.
\end{itemize}
}
}
\vspace{0.5cm}
\fbox{
\parbox{12cm}{
\subsection*{$v_i$}
\subsubsection*{\underline{Element}}
\texttt{<lpos:vi>}
\subsubsection*{\underline{Attributes}}
\texttt{mode}: a positive integer, identifying the normal mode that this quantum number is associated with.
\subsubsection*{\underline{Description}}
$v_i$ is the vibrational quantum number associated with the $\nu_i$ normal mode.
\subsubsection*{\underline{Restrictions}}
\begin{itemize}
\item non-negative integer.
\end{itemize}
}
}
\vspace{0.5cm}
\fbox{
\parbox{12cm}{
\subsection*{$l_i$}
\subsubsection*{\underline{Element}}
\texttt{<lpos:li>}
\subsubsection*{\underline{Attributes}}
\texttt{mode}: a positive integer, identifying the degenerate normal mode that this vibrational angular momentum quantum number is associated with.
\subsubsection*{\underline{Description}}
$l_i$ is the vibrational angular momentum quantum number associated with the degenerate $\nu_i$ normal mode; positive and negative values distinguish $l$-type doubling components; if two or more degenerate vibrations are excited, $l_i$ is only approximately defined (\emph{i.e.} it is not a totally good quantum number) - see \emph{e.g.} Herzerg II, p.212.
\subsubsection*{\underline{Restrictions}}
\begin{itemize}
\item non-negative integer.
\item $|l_i| = v_i, v_i-2, \cdots, 1 \;\mathrm{or}\;0$.
\end{itemize}
}
}
\vspace{0.5cm}
\fbox{
\parbox{12cm}{
\subsection*{$l$}
\subsubsection*{\underline{Element}}
\texttt{<lpos:l>}
\subsubsection*{\underline{Attributes}}
None.
\subsubsection*{\underline{Description}}
$l$ is the total vibrational angular momentum quantum number associated with resultant vibrational angular momentum about the internuclear axis.
\subsubsection*{\underline{Restrictions}}
\begin{itemize}
\item non-negative integer.
\end{itemize}
}
}
\vspace{0.5cm}
\fbox{
\parbox{12cm}{
\subsection*{\texttt{vibInv}}
\subsubsection*{\underline{Element}}
\texttt{<lpos:vibInv>}
\subsubsection*{\underline{Attributes}}
None.
\subsubsection*{\underline{Description}}
\texttt{vibInv} is the parity of the vibronic wavefunction with respect to inversion through the molecular centre of mass in the molecular coordinate system.
\subsubsection*{\underline{Restrictions}}
\begin{itemize}
\item `g' or `u'.
\end{itemize}
}
}
\vspace{0.5cm}
\fbox{
\parbox{12cm}{
\subsection*{\texttt{vibRefl}}
\subsubsection*{\underline{Element}}
\texttt{<lpos:vibRefl>}
\subsubsection*{\underline{Attributes}}
None.
\subsubsection*{\underline{Description}}
\texttt{vibRefl} is the parity of the vibronic wavefunction with respect to reflection in a plane containing the molecular symmetry axis in the molecular coordinate system.
\subsubsection*{\underline{Restrictions}}
\begin{itemize}
\item `$+$' or `$-$'.
\end{itemize}
}
}
\vspace{0.5cm}
\fbox{
\parbox{12cm}{
\subsection*{$J$}
\subsubsection*{\underline{Element}}
\texttt{<lpos:J>}
\subsubsection*{\underline{Attributes}}
None.
\subsubsection*{\underline{Description}}
$J$ is the quantum number associated with the total angular momentum excluding nuclear spin, $\boldsymbol{J}$.
\subsubsection*{\underline{Restrictions}}
\begin{itemize}
\item non-negative integer or half-odd integer.
\item $J \ge |l|$.
\end{itemize}
}
}
\vspace{0.5cm}
\fbox{
\parbox{12cm}{
\subsection*{$N$}
\subsubsection*{\underline{Element}}
\texttt{<lpos:N>}
\subsubsection*{\underline{Attributes}}
None.
\subsubsection*{\underline{Description}}
$N$ is the quantum number associated with the total angular momentum excluding electronic and nuclear spin, N: $\boldsymbol{J} = \boldsymbol{N} + \boldsymbol{S}$.
\subsubsection*{\underline{Restrictions}}
\begin{itemize}
\item non-negative integer.
\end{itemize}
}
}
\vspace{0.5cm}
\fbox{
\parbox{12cm}{
\subsection*{$I$}
\subsubsection*{\underline{Element}}
\texttt{<lpos:I>}
\subsubsection*{\underline{Attributes}}
\texttt{nuclearSpinRef}: a label, matching \texttt{/Q.+/} identifying the group of nuclear spins coupled to one another to form a total nuclear spin angular momentum.
\subsubsection*{\underline{Description}}
$I$ is the quantum number associated with the total nuclear spin angular momentum: $\boldsymbol{I} = \boldsymbol{I_1} + \boldsymbol{I_2} + \cdots$ where nuclei $1, 2, \cdots$ have individual nuclear spin angular momenta $\boldsymbol{I_1}, \boldsymbol{I_2}, \cdots$.
\subsubsection*{\underline{Restrictions}}
\begin{itemize}
\item non-negative integer or half-integer.
\end{itemize}
}
}
\vspace{0.5cm}
\fbox{
\parbox{12cm}{
\subsection*{$F_j$}
\subsubsection*{\underline{Element}}
\texttt{<lpos:Fj>}
\subsubsection*{\underline{Attributes}}
\begin{itemize}\item\texttt{nuclearSpinRef}: a label identifying the nuclear spin being coupled to $\boldsymbol{J}$ or $\boldsymbol{F_{j-1}}$ to form an intermediate angular momentum;\item $j$: an integer label identifying the order of the hyperfine coupling\end{itemize}.
\subsubsection*{\underline{Description}}
$F_j$ is the intermediate angular momentum quantum number associated with the coupling of the nuclear spin angular momentum of nucleus $j$ to the intermediate angular momentum: $\boldsymbol{F_1} = \boldsymbol{J} + \boldsymbol{I_1}$ or $\boldsymbol{F_j} = \boldsymbol{F_{j-1}} + \boldsymbol{I_j}$; $\boldsymbol{F_j}$ is often not a good quantum number.
\subsubsection*{\underline{Restrictions}}
\begin{itemize}
\item non-negative integer or half-integer.
\end{itemize}
}
}
\vspace{0.5cm}
\fbox{
\parbox{12cm}{
\subsection*{$F$}
\subsubsection*{\underline{Element}}
\texttt{<lpos:F>}
\subsubsection*{\underline{Attributes}}
\texttt{nuclearSpinRef}: a label, matching \texttt{/Q.+/} identifying the nuclear spin being coupled to $\boldsymbol{J}$ or $\boldsymbol{F_j}$ to form the total angular momentum.
\subsubsection*{\underline{Description}}
$F$ is the quantum number associated with the total angular momentum including nuclear spin: $\boldsymbol{F} = \boldsymbol{J} + \boldsymbol{I_1}$ if only one such coupling is resolved, $\boldsymbol{F} = \boldsymbol{F_{j-1}} + \boldsymbol{I_j}$ if two or more such couplings are resolved.
\subsubsection*{\underline{Restrictions}}
\begin{itemize}
\item non-negative integer or half-integer.
\end{itemize}
}
}
\vspace{0.5cm}
\fbox{
\parbox{12cm}{
\subsection*{$r$}
\subsubsection*{\underline{Element}}
\texttt{<lpos:r>}
\subsubsection*{\underline{Attributes}}
\texttt{name}: a string identifying this ranking index.
\subsubsection*{\underline{Description}}
$r$ is a named, positive integer label identifying the state if no other good quantum numbers or symmetries are known.
\subsubsection*{\underline{Restrictions}}
\begin{itemize}
\item positive integer.
\end{itemize}
}
}
\vspace{0.5cm}
\fbox{
\parbox{12cm}{
\subsection*{\texttt{parity}}
\subsubsection*{\underline{Element}}
\texttt{<lpos:parity>}
\subsubsection*{\underline{Attributes}}
None.
\subsubsection*{\underline{Description}}
\texttt{parity} is the total parity: the parity of the total molecular wavefunction (excluding nuclear spin) with respect to inversion through the molecular centre of mass of all particles' coordinates in the laboratory coordinate system, the $\hat{E}^*$ operation.
\subsubsection*{\underline{Restrictions}}
\begin{itemize}
\item `$+$' or `$-$'.
\end{itemize}
}
}
\vspace{0.5cm}
\fbox{
\parbox{12cm}{
\subsection*{\texttt{kronigParity}}
\subsubsection*{\underline{Element}}
\texttt{<lpos:kronigParity>}
\subsubsection*{\underline{Attributes}}
None.
\subsubsection*{\underline{Description}}
\texttt{kronigParity} is the `rotationless' parity: the parity of the total molecular wavefunction excluding nuclear spin and rotation with respect to inversion through the molecular centre of mass of all particles' coordinates in the laboratory coordinate system. For integral $J$, the levels are called: \begin{align*}e\;\mathrm{if\;parity\;is}\;+(-1)^J,\\f\;\mathrm{if\;parity\;is}\;-(-1)^J.\end{align*}For half-odd integer $J$, the levels are called: \begin{align*}e\;\mathrm{if\;parity\;is}\;+(-1)^{J-\frac{1}{2}},\\f\;\mathrm{if\;parity\;is}\;-(-1)^{J-\frac{1}{2}}.\end{align*}.
\subsubsection*{\underline{Restrictions}}
\begin{itemize}
\item `e' or `f'.
\end{itemize}
}
}
\vspace{0.5cm}
\fbox{
\parbox{12cm}{
\subsection*{\texttt{asSym}}
\subsubsection*{\underline{Element}}
\texttt{<lpos:asSym>}
\subsubsection*{\underline{Attributes}}
None.
\subsubsection*{\underline{Description}}
\texttt{asSym} is (for linear molecules with a centre of inversion) the symmetry of the rovibronic wavefunction: `a' or `s' such that the total wavefunction including nuclear spin is symmetric or antisymmetric with respect to permutation of the identical nuclei ($\hat{P}_{12}$), according to whether they are bosons or fermions respectively.
\subsubsection*{\underline{Restrictions}}
\begin{itemize}
\item `s' or `a'.
\end{itemize}
}
}
\vspace{0.5cm}




\section*{Open-shell, non-linear triatomics: nltos}
\fbox{
\parbox{12cm}{
\subsection*{\texttt{ElecStateLabel}}
\subsubsection*{\underline{Element}}
\texttt{<nltos:ElecStateLabel>}
\subsubsection*{\underline{Attributes}}
None.
\subsubsection*{\underline{Description}}
\texttt{ElecStateLabel} is a label identifying the electronic state: $X$, $A$, $a$, $B$, etc..
\subsubsection*{\underline{Restrictions}}
\begin{itemize}
\item string.
\end{itemize}
}
}
\vspace{0.5cm}
\fbox{
\parbox{12cm}{
\subsection*{\texttt{elecSym}}
\subsubsection*{\underline{Element}}
\texttt{<nltos:elecSym>}
\subsubsection*{\underline{Attributes}}
\texttt{group}: the symmetry group to which this species belongs.
\subsubsection*{\underline{Description}}
\texttt{elecSym} is the symmetry species of the electronic wavefunction described by some symmetry group..
\subsubsection*{\underline{Restrictions}}
\begin{itemize}
\item string.
\end{itemize}
}
}
\vspace{0.5cm}
\fbox{
\parbox{12cm}{
\subsection*{\texttt{elecInv}}
\subsubsection*{\underline{Element}}
\texttt{<nltos:elecInv>}
\subsubsection*{\underline{Attributes}}
None.
\subsubsection*{\underline{Description}}
\texttt{elecInv} is the parity of the electronic wavefunction with respect to inversion through the molecular centre of mass in the molecular coordinate system..
\subsubsection*{\underline{Restrictions}}
\begin{itemize}
\item `g' or `u'.
\end{itemize}
}
}
\vspace{0.5cm}
\fbox{
\parbox{12cm}{
\subsection*{$S$}
\subsubsection*{\underline{Element}}
\texttt{<nltos:S>}
\subsubsection*{\underline{Attributes}}
None.
\subsubsection*{\underline{Description}}
$S$ is the quantum number associated with the total electronic spin angular momentum, $\boldsymbol{S}$.
\subsubsection*{\underline{Restrictions}}
\begin{itemize}
\item non-negative integer or half-odd integer.
\end{itemize}
}
}
\vspace{0.5cm}
\fbox{
\parbox{12cm}{
\subsection*{$v_1$}
\subsubsection*{\underline{Element}}
\texttt{<nltos:v1>}
\subsubsection*{\underline{Attributes}}
None.
\subsubsection*{\underline{Description}}
$v_1$ is the vibrational quantum number associated with the $\nu_1$ normal mode.
\subsubsection*{\underline{Restrictions}}
\begin{itemize}
\item non-negative integer.
\end{itemize}
}
}
\vspace{0.5cm}
\fbox{
\parbox{12cm}{
\subsection*{$v_2$}
\subsubsection*{\underline{Element}}
\texttt{<nltos:v2>}
\subsubsection*{\underline{Attributes}}
None.
\subsubsection*{\underline{Description}}
$v_2$ is the vibrational quantum number associated with the $\nu_2$ normal mode.
\subsubsection*{\underline{Restrictions}}
\begin{itemize}
\item non-negative integer.
\end{itemize}
}
}
\vspace{0.5cm}
\fbox{
\parbox{12cm}{
\subsection*{$v_3$}
\subsubsection*{\underline{Element}}
\texttt{<nltos:v3>}
\subsubsection*{\underline{Attributes}}
None.
\subsubsection*{\underline{Description}}
$v_3$ is the vibrational quantum number associated with the $\nu_3$ normal mode.
\subsubsection*{\underline{Restrictions}}
\begin{itemize}
\item non-negative integer.
\end{itemize}
}
}
\vspace{0.5cm}
\fbox{
\parbox{12cm}{
\subsection*{$J$}
\subsubsection*{\underline{Element}}
\texttt{<nltos:J>}
\subsubsection*{\underline{Attributes}}
None.
\subsubsection*{\underline{Description}}
$J$ is the quantum number associated with the total angular momentum excluding nuclear spin, $\boldsymbol{J}$.
\subsubsection*{\underline{Restrictions}}
\begin{itemize}
\item non-negative integer or half-odd integer.
\end{itemize}
}
}
\vspace{0.5cm}
\fbox{
\parbox{12cm}{
\subsection*{$N$}
\subsubsection*{\underline{Element}}
\texttt{<nltos:N>}
\subsubsection*{\underline{Attributes}}
None.
\subsubsection*{\underline{Description}}
$N$ is the quantum number associated with the total angular momentum excluding electronic and nuclear spin, N: $\boldsymbol{J} = \boldsymbol{N} + \boldsymbol{S}$.
\subsubsection*{\underline{Restrictions}}
\begin{itemize}
\item non-negative integer.
\end{itemize}
}
}
\vspace{0.5cm}
\fbox{
\parbox{12cm}{
\subsection*{$K_a$}
\subsubsection*{\underline{Element}}
\texttt{<nltos:Ka>}
\subsubsection*{\underline{Attributes}}
None.
\subsubsection*{\underline{Description}}
$K_a$ is the rotational quantum label of an asymmetric top molecule, correlating to $K$ in the prolate symmetric top limit.
\subsubsection*{\underline{Restrictions}}
\begin{itemize}
\item non-negative integer.
\item $K_a \le J$.
\item $K_a + K_c = J \;\mathrm{or}\;J + 1$.
\end{itemize}
}
}
\vspace{0.5cm}
\fbox{
\parbox{12cm}{
\subsection*{$K_c$}
\subsubsection*{\underline{Element}}
\texttt{<nltos:Kc>}
\subsubsection*{\underline{Attributes}}
None.
\subsubsection*{\underline{Description}}
$K_c$ is the rotational quantum label of an asymmetric top molecule, correlating to $K$ in the oblate symmetric top limit.
\subsubsection*{\underline{Restrictions}}
\begin{itemize}
\item non-negative integer.
\item $K_c \le J$.
\item $K_a + K_c = J \;\mathrm{or}\;J + 1$.
\end{itemize}
}
}
\vspace{0.5cm}
\fbox{
\parbox{12cm}{
\subsection*{$F_1$}
\subsubsection*{\underline{Element}}
\texttt{<nltos:F1>}
\subsubsection*{\underline{Attributes}}
\texttt{nuclearSpinRef}: a label identifying the nuclear spin coupled to $\boldsymbol{J}$ to form the intermediate angular momentum.
\subsubsection*{\underline{Description}}
$F_1$ is the intermediate angular momentum quantum number associated with the coupling of the rotational angular momentum and nuclear spin of nucleus 1 where two or more such couplings are resolved: $\boldsymbol{F_1} = \boldsymbol{J} + \boldsymbol{I_1}$; $F_1$ is often not a good quantum number.
\subsubsection*{\underline{Restrictions}}
\begin{itemize}
\item non-negative integer or half-integer.
\item $|J - I_1| \le F_1 \le J + I_1$.
\end{itemize}
}
}
\vspace{0.5cm}
\fbox{
\parbox{12cm}{
\subsection*{$F_2$}
\subsubsection*{\underline{Element}}
\texttt{<nltos:F2>}
\subsubsection*{\underline{Attributes}}
\texttt{nuclearSpinRef}: a label identifying the nuclear spin coupled to $\boldsymbol{F_1}$ to form an intermediate angular momentum.
\subsubsection*{\underline{Description}}
$F_2$ is the intermediate angular momentum quantum number associated with the coupling of the rotational angular momentum and nuclear spin of nucleus 2 where three such couplings are resolved: $\boldsymbol{F_2} = \boldsymbol{F_1} + \boldsymbol{I_2}$; $F_2$ is often not a good quantum number.
\subsubsection*{\underline{Restrictions}}
\begin{itemize}
\item non-negative integer or half-integer.
\item $|F_1 - I_2| \le F_2 \le F_1 + I_2$.
\end{itemize}
}
}
\vspace{0.5cm}
\fbox{
\parbox{12cm}{
\subsection*{$F$}
\subsubsection*{\underline{Element}}
\texttt{<nltos:F>}
\subsubsection*{\underline{Attributes}}
\texttt{nuclearSpinRef}: a label identifying the nuclear spin coupled to $\boldsymbol{J}$, $\boldsymbol{F_1}$, or $\boldsymbol{F}$ to form the total angular momentum.
\subsubsection*{\underline{Description}}
$F$ is the quantum number associated with the total angular momentum including nuclear spin: $\boldsymbol{F} = \boldsymbol{J} + \boldsymbol{I_1}$ if only one hyperfine coupling is resolved, $\boldsymbol{F} = \boldsymbol{F_1} + \boldsymbol{I_2}$ if two couplings are resolved, or $\boldsymbol{F} = \boldsymbol{F_2} + \boldsymbol{I_3}$ if all three couplings are resolved.
\subsubsection*{\underline{Restrictions}}
\begin{itemize}
\item non-negative integer or half-integer.
\item $|F_2 - I_3| \le F \le F_2 + I_3$.
\end{itemize}
}
}
\vspace{0.5cm}
\fbox{
\parbox{12cm}{
\subsection*{$r$}
\subsubsection*{\underline{Element}}
\texttt{<nltos:r>}
\subsubsection*{\underline{Attributes}}
\texttt{name}: a string identifying this ranking index.
\subsubsection*{\underline{Description}}
$r$ is a named, positive integer label identifying the state if no other good quantum numbers or symmetries are known.
\subsubsection*{\underline{Restrictions}}
\begin{itemize}
\item positive integer.
\end{itemize}
}
}
\vspace{0.5cm}
\fbox{
\parbox{12cm}{
\subsection*{\texttt{parity}}
\subsubsection*{\underline{Element}}
\texttt{<nltos:parity>}
\subsubsection*{\underline{Attributes}}
None.
\subsubsection*{\underline{Description}}
\texttt{parity} is the total parity: the parity of the total molecular wavefunction (excluding nuclear spin) with respect to inversion through the molecular centre of mass of all particles' coordinates in the laboratory coordinate system, the $\hat{E}^*$ operation.
\subsubsection*{\underline{Restrictions}}
\begin{itemize}
\item `$+$' or `$-$'.
\end{itemize}
}
}
\vspace{0.5cm}
\fbox{
\parbox{12cm}{
\subsection*{\texttt{asSym}}
\subsubsection*{\underline{Element}}
\texttt{<nltos:asSym>}
\subsubsection*{\underline{Attributes}}
None.
\subsubsection*{\underline{Description}}
\texttt{asSym} is the symmetry of the rovibronic wavefunction: `a' or `s' such that the total wavefunction including nuclear spin is symmetric or antisymmetric with respect to permutation of the identical nuclei ($\hat{P}_{12}$), according to whether they are bosons or fermions respectively.
\subsubsection*{\underline{Restrictions}}
\begin{itemize}
\item `s' or `a'.
\end{itemize}
}
}
\vspace{0.5cm}





\end{document}
